\documentclass[10pt,reqno]{amsart}

\usepackage[T1]{fontenc}
\usepackage{amsmath}
\usepackage{amssymb}
\usepackage{mathtools}
\usepackage{booktabs}
\usepackage{array}
\usepackage{enumitem}
\usepackage{hyperref}
\usepackage{url}
\def\UrlBreaks{\do\_\do\/\do\-}

\hypersetup{
  hidelinks,
  pdftitle={Identity Elimination with Observable Metadata, and the Collapse of a Computational-Path Rewrite Quotient},
  pdfauthor={Arthur Freitas Ramos, Ruy J. G. B. de Queiroz, Anjolina Grisi de Oliveira, and Dov M. Gabbay}
}
\allowdisplaybreaks
\emergencystretch=3em
\setlength{\abovedisplayskip}{7pt plus 2pt minus 3pt}
\setlength{\belowdisplayskip}{7pt plus 2pt minus 3pt}
\setlength{\abovedisplayshortskip}{4pt plus 2pt}
\setlength{\belowdisplayshortskip}{5pt plus 2pt minus 2pt}

\title[Metadata, Identity Elimination, and a Quotient Collapse]
  {Identity Elimination with Observable Metadata,\\
   and the Collapse of a Computational-Path Rewrite Quotient}

\author{Arthur Freitas Ramos}
\author{Ruy J. G. B. de Queiroz}
\author{Anjolina Grisi de Oliveira}
\author{Dov M. Gabbay}
\renewcommand{\shortauthors}{Ramos et al.}

\subjclass[2020]{03B15, 03B38, 03B70, 18N45}
\keywords{identity elimination, path induction, observable metadata,
setoid quotient, contractibility, axiom K, computational paths, equality
traces}

\newtheorem{theorem}{Theorem}[section]
\newtheorem{proposition}[theorem]{Proposition}
\newtheorem{lemma}[theorem]{Lemma}
\newtheorem{corollary}[theorem]{Corollary}
\newtheorem{definition}[theorem]{Definition}
\newtheorem{example}[theorem]{Example}
\newtheorem{remark}[theorem]{Remark}
\newtheorem{warning}[theorem]{Warning}

\newcommand{\U}{\mathcal U}
\newcommand{\V}{\mathcal V}
\newcommand{\isContr}{\mathsf{isContr}}
\newcommand{\Contr}{\mathsf{Contr}}
\newcommand{\Eqv}{\simeq}
\newcommand{\Jbeta}{\mathsf J^\beta}
\newcommand{\Meta}{\mathsf M}
\newcommand{\Tot}{\mathsf T}
\newcommand{\RwTot}{\mathsf{RwTot}}
\newcommand{\Step}{\mathsf{Step}}
\newcommand{\PathTy}{\mathsf{Path}}
\newcommand{\refl}{\mathsf{refl}}
\newcommand{\transport}{\mathsf{transport}}
\newcommand{\List}{\mathsf{List}}
\newcommand{\proj}{\pi_{=}}
\newcommand{\Unit}{\mathbf 1}
\newcommand{\Empty}{[\,]}
\newcommand{\SetTot}{\mathsf{SetoidTotal}}
\newcommand{\Ker}{\mathsf{Ker}}
\newcommand{\PRW}{\mathsf{PathRwQuot}}
\newcommand{\LocalK}{\mathsf K_{\mathrm{loc}}}
\newcommand{\Q}{\mathsf Q}

\begin{document}

\begin{abstract}
Suppose that equality from a fixed point \(a:A\) to \(b:A\) is refined by
observable metadata:
\[
  R_M(a,b)=\sum_{h:a=b}M(b,h),
  \qquad r^M_a=(\refl_a,m_0).
\]
We first diagnose exactly when this refinement supports unrestricted based
elimination into arbitrary motives with a propositional beta law:
\[
 \Jbeta_{R_M}(a,r^M_a)
 \quad\Longleftrightarrow\quad
 \isContr\bigl(M(a,\refl_a)\bigr).
\]
The underlying pointed-family equivalence is the proof-irrelevant,
propositional-beta specialization of the identity-system characterization of
\cite[Thm.~5.8.2]{HoTTBook2013}; the contribution begins with the computation of
the decorated total space and its exact quotient classification.

We then classify every ordinary setoid quotient repair.  If \(S_{b,h}\) is a
setoid on \(M(b,h)\) and \(m_0:M(a,\refl_a)\), then the quotient-repaired
refinement supports the same eliminator exactly when the reflexivity relation
\(S_{a,\refl_a}\) is total.  For an arbitrary carrier \(X\), the precise
empty-sensitive statement is
\[
 \isContr(X/S)\quad\Longleftrightarrow\quad
 \bigl(\mathsf{Inh}(X)\ \wedge\ \SetTot(S)\bigr).
\]
Thus every set-level repair supporting unrestricted \(J\) identifies all
reflexivity metadata; no visible distinction can survive there.  The
classification extends beyond sigma-shaped refinements: for a reflexive
projection \(q_b:R(b)\to(a=b)\), unrestricted elimination is equivalent to
contractibility of the reflexivity kernel
\(\Ker(q)=\{r:R(a)\mid q_a(r)=\refl_a\}\).  Applied to
\(\PRW\), this yields the exact local criterion---elimination exists if and only
if the genuine loop quotient is trivial---and, pushed back to representatives, a
quotient-free form: it holds exactly when \(\mathsf{RwEq}\) relates every pair of
raw loops at the base point.

We then settle that condition for a formalized computational-path rewriting
system, and the answer is a collapse.  One primitive rule decides it: because
packaging pointwise paths into a path between functions records the \emph{empty}
trace, the function-application beta rule exhibits the trace-free path as a
one-step predecessor of an arbitrary path with the same endpoints.  Consequently
\(\mathsf{RwEq}\) is total on every fiber and
\[
 \PRW_A(a,b)\simeq(a=b),
\]
so the rewrite quotient retains no computational-path information whatever and
every invariant on it is constant.  It does support unrestricted based path
induction on every carrier---but precisely because it has become ambient
equality, which is the price the classification predicts for an ordinary
setoid repair.  A design theorem makes the trade-off general: a setoid is total
exactly when every invariant of it is constant, so for an inhabited carrier the
quotient is contractible exactly when it retains nothing and the two goals are
strictly opposed.  Sharpness is proved by a parity invariant: the groupoid
fragment of the rewriting system preserves trace-length parity and therefore
does not collapse, so the degeneracy is caused by the function-space rules
rather than by path algebra.

Finally we audit the computational-path examples.  The reflexivity fiber of the
trace record is computed exactly---it is \(\List(A)\)---so failure for raw
records is unconditional on every pointed carrier, in sharp contrast with the
quotient.  The current ``circle'' is a one-constructor inductive type and the current
``torus'' its product; their genuine \(\PRW\) loop quotients are contractible,
while the separately defined winding-expression quotients are equivalent to
\(\mathbb Z\) and \(\mathbb Z^2\).  The collapse theorem upgrades the resulting
no-bridge statement from those two carriers to \emph{all} of them.  A Lean
development verifies every result using only propositional extensionality and
quotient soundness.  The scope remains proof-irrelevant ambient equality, not
full Martin-L\"of or homotopy type theory.
\end{abstract}

\maketitle

\section{Introduction}
\label{sec:introduction}

Equality evidence is often decorated.  A proof assistant may retain a
derivation history; an optimizer may attach a cost; a checker may keep a
certificate; a normalizer may record a log; and a rewriting calculus may expose
the sequence of elementary steps.  The decoration is useful precisely because
it is observable.  That same observability changes the elimination principle.

Fix \(a:A\).  Given a family
\[
        M:\prod_{b:A}(a=b)\longrightarrow\U
\]
and a chosen \(m_0:M(a,\refl_a)\), define
\[
 R_M(a,b)=\sum_{h:a=b}M(b,h),
 \qquad r^M_a=(\refl_a,m_0).
\tag{1.1}\label{eq:refinement}
\]
The projection
\[
       \proj:R_M(a,b)\longrightarrow(a=b),\qquad
       \proj(h,m)=h
\tag{1.2}\label{eq:projection}
\]
forgets the metadata.  Ordinary equality elimination can act through
\(\proj\).  The stronger question is whether \emph{every} motive on the
refined object can be eliminated from the single base case \(r^M_a\).

\begin{theorem}[Metadata-fiber criterion]
\label{thm:headline}
The refinement \(R_M\), based at \(r^M_a\), admits unrestricted based
elimination into all \(\U\)-valued motives with propositional beta if and only
if
\[
                        \isContr\bigl(M(a,\refl_a)\bigr).
\tag{1.3}\label{eq:headline}
\]
\end{theorem}

Theorem~\ref{thm:headline} organizes Part~\ref{part:general}.  Its two
directions are constructive here: an eliminator contracts the based total space
by choosing the motive ``the base point equals the current point'', and a
contraction transports the base datum to every point.  The classification step
is that the total space of \eqref{eq:refinement} is equivalent to the single
reflexivity metadata fiber, so all endpoint and equality coordinates disappear
while observable metadata over reflexivity remains.  This is what the familiar
contraction of \(\sum_{b:A}(a=b)\) does \emph{not} give: that argument explains
ordinary based path induction, but says nothing about a larger space obtained by
retaining an extra coordinate.  Theorem~\ref{thm:headline} measures exactly that
coordinate.

Our case study is the record
\[
 \PathTy_A(a,b)
   =\bigl\{\,L:\List(\Step(A)),\quad h:a=b\,\bigr\}
 \ \Eqv\ \sum_{h:a=b}\List(\Step(A)),
\tag{1.4}\label{eq:path-record}
\]
where \(L\) is a visible trace and the equivalence merely reorders fields.  At
reflexivity the empty list and the singleton step
\([\,a\xrightarrow{\refl_a}a\,]\) differ, so the fiber is not contractible and
no unrestricted eliminator can exist for motives observing \(L\).  The singleton
is a valid chain from \(a\) to \(a\), so imposing adjacency and endpoint
conditions does not help.

Two remedies follow: restrict motives to those factoring through \(\proj\),
which cannot observe the metadata and inherit equality elimination; or quotient
the metadata until its reflexivity quotient is contractible.  The first
preserves the data and limits observation, the second changes the identity-like
object.

The second option has a sharp universal boundary.  Write
\(\Q_S(X)=X/S\) and
\[
 \SetTot(S)\quad:\!\!\Longleftrightarrow\quad
 \prod_{x,y:X}S(x,y).
\tag{1.5}\label{eq:setoid-total-intro}
\]
Quotient soundness proves that totality contracts \(\Q_S(X)\), while quotient
exactness recovers \(S(x,y)\) from equality of the quotient images.  Hence, for
a chosen \(x_0:X\),
\[
 \isContr(\Q_S(X))\quad\Longleftrightarrow\quad\SetTot(S).
\tag{1.6}\label{eq:quotient-intro}
\]
Without \(x_0\), inhabitance must appear as an additional conjunct.  Applied to
the reflexivity metadata fiber, this is a classification of \emph{all} ordinary
setoid repairs, not only a sufficient construction.

We also replace the sigma-specific viewpoint by a projection theorem.  A
reflexive map \(q_b:R(b)\to(a=b)\) determines
\[
 \Ker(q)=\sum_{r:R(a)}\bigl(q_a(r)=\refl_a\bigr),
\tag{1.7}\label{eq:kernel-intro}
\]
and equality induction gives
\(\sum_bR(b)\Eqv\Ker(q)\).  Consequently unrestricted elimination is
equivalent to contractibility of this kernel.  The \(\PRW\) application is
especially useful: its kernel is equivalent to the entire loop quotient, so
the exact condition is local quotient-level axiom K, and, pushed back to
representatives, the quotient-free condition that \(\mathsf{RwEq}\) relate every
pair of raw loops.

That last condition is decidable by inspection of the rule set, and
Section~\ref{sec:collapse} decides it.  The computational-path implementation
packages a family of pointwise paths into a path between functions while
recording the \emph{empty} trace; the primitive application rule then makes the
trace-free path a one-step predecessor of an \emph{arbitrary} path with the same
endpoints.  Consequently \(\mathsf{RwEq}\) is total on every fiber and
\[
 \PRW_A(a,b)\Eqv(a=b),
\tag{1.8}\label{eq:collapse-intro}
\]
The rewrite quotient is therefore ambient equality.  It retains no
computational-path information, every invariant on it is constant, and it
supports unrestricted based path induction on every carrier for exactly that
reason.  We state the positive half with its cause attached throughout---the
implemented rewrite quotient supports path induction \emph{because} it is
equivalent to ambient equality---since the unqualified form invites the reading
that a proof-relevant path object has been shown to admit \(J\), which is not
what is proved.

Section~\ref{sec:design} shows the trade-off is structural rather than
accidental: a setoid is total exactly when every invariant of it is constant, so
for an inhabited carrier---which a pointed fiber always is---elimination and
retained information are strictly opposed, and any rule making one path a
predecessor of a whole fiber destroys everything there.  Section~\ref{subsec:sharpness} then localises the present failure: the
groupoid fragment preserves trace-length parity and does not collapse, so the
degeneracy comes from the function-space rules and not from path algebra.

\paragraph{Contributions.}
The centrepiece is Part~\ref{part:case-study}: the implemented rewrite quotient
collapses to ambient equality, so it retains no computational-path information
and admits only constant invariants
(Theorem~\ref{thm:rweq-total}, Corollary~\ref{cor:unconditional-repair}).  A
parity invariant shows the groupoid fragment does \emph{not} collapse, locating
the cause in the function-space rules (Theorem~\ref{thm:sharpness}), and a design
theorem shows the trade-off is structural: a quotient is contractible exactly
when every invariant of it is constant (Theorem~\ref{thm:design}).  The audit of
the circle and torus, and the universal no-bridge theorem
(Theorem~\ref{thm:universal-no-bridge}), follow.

Part~\ref{part:general} supplies the classification these rest on: the
total-space computation \(\sum_b\sum_hM(b,h)\Eqv M(a,\refl_a)\) and its
metadata-fiber criterion; the empty-sensitive quotient theorem
\(\isContr(X/S)\leftrightarrow(\mathsf{Inh}(X)\wedge\SetTot(S))\) with its
specialization to every setoid repair of metadata, the indiscrete repair, and a
worked nontrivial quotient that repairs nothing; the projection/kernel theorem
and its raw-level form, which makes the criterion checkable without forming any
quotient; and a factorization theorem for equality-projection motives.  These
ingredients are individually elementary; their value is that together they
decide the case study.

Theorem~\ref{thm:j-total}, the equivalence between unrestricted based
elimination and based total-space contractibility, is stated for completeness
but is not claimed as a contribution; see Remark~\ref{rem:standard-stage}.

\paragraph{Scope.}
We reason about proof-irrelevant ambient equality in an extensional
set/type-theoretic metatheory, matching the formal artifact.  We classify
\emph{propositional} beta only.  We do not present a full syntax or model of
Martin-L\"of type theory, do not assume or prove univalence or higher
inductive types, and do not identify \eqref{eq:path-record} with the HoTT
identity type.  The collapse theorem of Section~\ref{sec:collapse} is a
statement about the formalized rule set, not about computational-path calculi in
general; the raw de Bruijn calculus developed alongside this work is published
separately as a self-contained article \cite{ScopedTraceCalculus2026}, and
Section~\ref{sec:formalization} describes the division.

\paragraph{Organization.}
The paper is in two parts.  Part~\ref{part:general} is the general theory and
mentions no rewrite system: Section~\ref{sec:setting} fixes the metatheory,
Section~\ref{sec:total-space} recalls the identity-system classification,
Section~\ref{sec:metadata} proves Theorem~\ref{thm:headline},
Section~\ref{sec:setoid-repair} classifies setoid quotient repairs,
Section~\ref{sec:projection-kernel} proves the projection/kernel theorem,
Section~\ref{sec:factorization} gives the factorizing-motive theorem, and
Section~\ref{sec:examples} classifies further metadata.

Part~\ref{part:case-study} is the case study.  Section~\ref{sec:traces} computes
the raw trace fiber, Section~\ref{sec:pathrwquot-criterion} specializes the
kernel theorem to the rewrite quotient and removes the quotient from the
criterion, Section~\ref{sec:collapse} proves the collapse and its sharpness,
Section~\ref{sec:design} derives the design constraint any replacement rule set
must meet, and Section~\ref{sec:genuine-synthetic} audits the circle and torus
definitions.  Sections~\ref{sec:comparison}--\ref{sec:conclusion} state
boundaries, related work, formal verification, and conclusions.

\part{Classification of equality with observable metadata}
\label{part:general}

Part~\ref{part:general} is independent of computational paths.  It classifies which decorations of an equality retain unrestricted based
elimination, and which quotients of a decoration restore it.  Nothing in
it mentions a rewrite system.


\section{Setting and definitions}
\label{sec:setting}

\subsection{Ambient equality and universes}

Fix a universe \(\U\) closed under dependent sums and products, equality,
finite lists, and the elementary constructions below.  A larger universe
\(\V\) contains collections quantifying over all \(\U\)-valued motives.
Thus an individual motive value and every equality proposition remain
\(\U\)-small, while the record of an operation accepting every such motive
may live in \(\V\).  No impredicative identification \(\U=\V\) is intended.

Ambient equality has elimination, transport, function extensionality, and
proof irrelevance.  For \(e:x=y\) and \(C:X\to\U\), write
\[
                  \transport_C(e,-):C(x)\longrightarrow C(y).
\]
Its reflexivity law is used propositionally:
\[
                  \transport_C(\refl_x,c)=c.
\tag{2.1}\label{eq:transport-beta}
\]
Proof irrelevance implies that any two witnesses of the same ambient equality
are equal.  This assumption is appropriate for a proof-irrelevant equality, which
lives in \(\mathsf{Prop}\).  It is also why the present theory must not be read
as a claim about arbitrary proof-relevant identity types.

An equivalence \(e:X\Eqv Y\) consists of maps
\[
 e:X\to Y,\qquad e^{-1}:Y\to X
\]
with specified inverse equations.  We use no univalence principle: an
equivalence is data, not an equality of types.

\begin{definition}[Contractibility]
For \(c:X\), define
\[
 \Contr_c(X)=\prod_{x:X}(c=x),\qquad
 \isContr(X)=\sum_{c:X}\Contr_c(X).
\tag{2.2}\label{eq:contractibility}
\]
\end{definition}

The orientation \(c=x\) is convenient for transport from the center.  If
\((c,\kappa):\isContr(X)\), then every two points are equal:
\[
             x=(\kappa x)^{-1}\cdot\kappa y=y.
\tag{2.3}\label{eq:all-equal}
\]
For a specified \(c_0:X\), the propositions
\(\Contr_{c_0}(X)\) and \(\isContr(X)\) are equivalent.  One direction chooses
\(c_0\); the other uses
\[
      c_0=(\kappa c_0)^{-1}\cdot\kappa x=x.
\tag{2.4}\label{eq:recenter}
\]

\subsection{Based relations and unrestricted elimination}

Let \(R:A\to\U\) be a relation row based at \(a\), and choose \(r_a:R(a)\).
Its based total space and center are
\[
             B_R=\sum_{b:A}R(b),\qquad c_R=(a,r_a).
\tag{2.5}\label{eq:based-total}
\]
Writing \(R(b)\) rather than \(R(a,b)\) keeps the fixed base point implicit.

\begin{definition}[Unrestricted based elimination with propositional beta]
\label{def:unrestricted}
The data \(\Jbeta_R(a,r_a)\) consist of an operation
\[
 \mathcal J:
   \prod_{C:B_R\to\U}
   C(c_R)\longrightarrow\prod_{x:B_R}C(x)
\tag{2.6}\label{eq:j-data}
\]
and equations
\[
       \beta_{C,d}:\mathcal J(C,d,c_R)=d
       \qquad(C:B_R\to\U,\ d:C(c_R)).
\tag{2.7}\label{eq:j-beta}
\]
\end{definition}

The adjective ``unrestricted'' refers only to the motive: \(C\) may inspect
the entire point \(x=(b,p)\), including every field of \(p\).  The beta law is
an ambient equality, not a judgmental reduction.  Definition
\ref{def:unrestricted} is based: \(a\) and \(r_a\) are fixed.

\section{Based elimination and total-space contractibility}
\label{sec:total-space}

\begin{theorem}[Total-space classification]
\label{thm:j-total}
For any based relation row \((R,r_a)\),
\[
        \Jbeta_R(a,r_a)
        \quad\Longleftrightarrow\quad
        \Contr_{c_R}(B_R)
        \quad\Longleftrightarrow\quad
        \isContr(B_R).
\tag{3.1}\label{eq:j-total}
\]
Here the first term means that the data of Definition
\ref{def:unrestricted} exist.
\end{theorem}

\begin{proof}
Assume \((\mathcal J,\beta)\).  Use the motive
\[
                         C(x)=(c_R=x).
\tag{3.2}\label{eq:contraction-motive}
\]
Its base datum is \(\refl_{c_R}:C(c_R)\).  Define
\[
       \kappa(x)=\mathcal J(C,\refl_{c_R},x):c_R=x.
\tag{3.3}\label{eq:kappa-from-j}
\]
This is an element of \(\Contr_{c_R}(B_R)\).  The supplied beta equation
also gives \(\kappa(c_R)=\refl_{c_R}\), although only the contraction itself
is needed for this direction.

Conversely, let \(\kappa:\Contr_{c_R}(B_R)\).  Given \(C\), \(d:C(c_R)\),
and \(x:B_R\), put
\[
        \mathcal J_\kappa(C,d,x)
        =\transport_C(\kappa(x),d):C(x).
\tag{3.4}\label{eq:j-from-kappa}
\]
At \(x=c_R\), proof irrelevance identifies
\(\kappa(c_R)\) with \(\refl_{c_R}\).  Congruence of transport in its equality
argument and~\eqref{eq:transport-beta} give
\[
 \mathcal J_\kappa(C,d,c_R)
 =\transport_C(\kappa(c_R),d)
 =\transport_C(\refl_{c_R},d)
 =d.
\tag{3.5}\label{eq:derived-beta}
\]
Thus \((\mathcal J_\kappa,\beta^\kappa)\) is the required data.
The final equivalence in~\eqref{eq:j-total} is the recentering argument
\eqref{eq:recenter}.
\end{proof}

\begin{remark}[Relation to identity systems]
\label{rem:standard-stage}
Theorem~\ref{thm:j-total} is not new, and we do not claim it as a contribution.
It is the proof-irrelevant, propositional-beta specialization of the standard
\emph{identity-system} characterization: a pointed family admits based induction
precisely when its total space is contractible.  In homotopy type theory this is
\cite[Thm.~5.8.2]{HoTTBook2013}, itself an instance of the initial-algebra
analysis of inductive identity types of Awodey, Gambino and Sojakova
\cite{AwodeyGambinoSojakova2012,AwodeyGambinoSojakova2017}; the underlying
universal-property argument for a pointed contractible space is older still.
Applied to \(R(b)=(a=b)\), equality elimination proves that
the based total identity space is contractible.  Neither fact mentions
metadata.

We record it in the present setting only because the ambient equality here is
proof-irrelevant and the beta law is propositional, so the standard statement
does not literally apply and its two directions become, respectively, a choice of
motive and a transport.  Our contribution begins in
Section~\ref{sec:metadata}, where a second equivalence \emph{computes} the
enlarged total space of a decorated equality, and continues through the exact
setoid-quotient classification of Section~\ref{sec:setoid-repair}, the
kernel form of Section~\ref{sec:projection-kernel}, and the concrete
\(\PRW\) analysis of Sections~\ref{sec:collapse}--\ref{sec:genuine-synthetic}.
Readers who take Theorem~\ref{thm:j-total} as known may begin at
Section~\ref{sec:metadata}.
\end{remark}

\begin{corollary}[Observable-point obstruction]
\label{cor:observable-point}
If \(x:B_R\) satisfies \(x\ne c_R\), then \(\Jbeta_R(a,r_a)\) does not exist.
More generally, it is enough to give \(O:B_R\to\mathsf{Prop}\) with
\(O(c_R)\) and \(\neg O(x)\).
\end{corollary}

\begin{proof}
An eliminator gives \(c_R=x\) by~\eqref{eq:kappa-from-j}, contradicting the
first hypothesis.  For the predicate formulation, transport \(O(c_R)\) along
that equality.  Equivalently, use \(O\) itself as the motive in
\eqref{eq:j-data}.
\end{proof}

The corollary is useful for finding counterexamples, but it does not yet
classify a structured family.  For equality refinements, the entire
obstruction can be localized to one fiber.

\section{The metadata-fiber criterion}
\label{sec:metadata}

Return to the data of~\eqref{eq:refinement}.  The based total space is
\[
       T_M(a)
       =\sum_{b:A}R_M(a,b)
       =\sum_{b:A}\sum_{h:a=b}M(b,h),
\tag{4.1}\label{eq:metadata-total}
\]
with center
\[
                         c_M=(a,\refl_a,m_0).
\tag{4.2}\label{eq:metadata-center}
\]

\subsection{Collapsing the equality coordinates}

\begin{lemma}[Based equality total space]
\label{lem:equality-total}
The type \(P_a=\sum_{b:A}(a=b)\) is contractible with center
\((a,\refl_a)\).
\end{lemma}

\begin{proof}
For \((b,h):P_a\), equality elimination on \(h\) reduces the goal
\[
                    (a,\refl_a)=(b,h)
\]
to reflexivity at \((a,\refl_a)\).  This defines the contraction for every
\((b,h)\).
\end{proof}

\begin{theorem}[Total space--fiber equivalence]
\label{thm:total-fiber}
There is a specified equivalence
\[
             \Phi_M:T_M(a)\Eqv M(a,\refl_a)
\tag{4.3}\label{eq:total-fiber}
\]
whose forward map sends \((a,\refl_a,m)\) to \(m\), and whose inverse sends
\(m\) to \((a,\refl_a,m)\).
\end{theorem}

\begin{proof}
The inverse is \(G(m)=(a,\refl_a,m)\).  For the forward map take
\((b,h,m):T_M(a)\) and eliminate on \(h:a=b\), returning \(m\) in the reflexive
case; equivalently
\(F(b,h,m)=\transport_{\widetilde M}(\gamma_{b,h}^{-1},m)\) using the
contraction \(\gamma_{b,h}\) of Lemma~\ref{lem:equality-total}.  Then
\(F(G(m))=m\) at reflexivity, and the other composite reduces to reflexivity
after eliminating on \(h\).
\end{proof}

The proof uses equality elimination only.  It chooses no representatives,
normalizes no metadata, and does not assume all metadata values equal \(m_0\):
it says every metadata value in the total space transports to the reflexivity
fiber, where it remains visible.

\begin{lemma}[Contractibility respects equivalence]
\label{lem:contr-equiv}
If \(e:X\Eqv Y\) then \(\isContr(X)\Longleftrightarrow\isContr(Y)\).
\end{lemma}

\begin{proof}
Transport the center and compose the contraction with the round-trip witness;
apply the same construction to \(e^{-1}\) for the converse.
\end{proof}

\begin{theorem}[Metadata-fiber criterion, full form]
\label{thm:metadata-fiber}
For \(R_M\) and \(r^M_a\) as in~\eqref{eq:refinement},
\[
 \Jbeta_{R_M}(a,r^M_a)
 \quad\Longleftrightarrow\quad
 \isContr\bigl(T_M(a)\bigr)
 \quad\Longleftrightarrow\quad
 \isContr\bigl(M(a,\refl_a)\bigr).
\tag{4.4}\label{eq:metadata-full}
\]
\end{theorem}

\begin{proof}
The first equivalence is Theorem~\ref{thm:j-total}, applied to the relation row
\(b\mapsto R_M(a,b)\) and center \(c_M\).  The second is
Lemma~\ref{lem:contr-equiv}, applied to the explicit equivalence
\(\Phi_M\) of Theorem~\ref{thm:total-fiber}.  Composing the equivalences proves
the claim.
\end{proof}

Theorem~\ref{thm:headline} is the first-to-last equivalence in
\eqref{eq:metadata-full}.  The chosen base metadata \(m_0\) does not by itself
need to equal every metadata value by definition; the criterion says exactly
that such equalities must be available as a contraction.

\subsection{Immediate classifications}

\begin{corollary}[Contractible decorations]
\label{cor:contractible-decoration}
If \(M(a,\refl_a)\) is contractible, then \(R_M\) admits unrestricted based
elimination with propositional beta.  In particular this holds when the
reflexivity fiber is:
\begin{enumerate}[label=(\roman*)]
\item the unit type;
\item an inhabited subsingleton;
\item a proof-irrelevant proposition with a witness; or
\item a type equipped with a quotient identifying every element with \(m_0\).
\end{enumerate}
\end{corollary}

\begin{corollary}[Visible alternatives]
\label{cor:visible-alternatives}
If \(m_1:M(a,\refl_a)\) and \(m_0\ne m_1\), unrestricted based elimination
with propositional beta does not exist.
\end{corollary}

\begin{proof}
A contraction would identify \(m_0\) and \(m_1\) by
\eqref{eq:all-equal}; apply Theorem~\ref{thm:metadata-fiber}.
\end{proof}

\begin{remark}[Only the reflexivity fiber]
\label{rem:only-fiber}
The criterion does not require a separate contractibility hypothesis for
every \(M(b,h)\).  Equality elimination transports the entire total space to
\(M(a,\refl_a)\).  Conversely, noncanonical metadata over reflexivity is
already fatal, even if every nonreflexive-looking presentation has a unique
decoration.
\end{remark}

\section{Universal setoid quotient repair}
\label{sec:setoid-repair}

The phrase ``quotient the metadata'' is too weak to be a theorem: many
relations can be quotiented, and a quotient can retain several reflexivity
classes.  We now classify exactly which ordinary setoid quotients repair
unrestricted elimination.  Quotient exactness is essential in the converse.

\subsection{The empty-sensitive quotient theorem}

\begin{definition}[Total setoid]
\label{def:setoid-total}
For a setoid \(S\) on \(X\), write
\(\SetTot(S):\Longleftrightarrow\prod_{x,y:X}S(x,y)\), and \(\eta_S:X\to X/S\)
for the canonical quotient map.
\end{definition}

Totality concerns the relation, not the number of raw terms: a type may have
many inhabitants while a total setoid has one class, and on the empty type the
condition holds vacuously.  That is why the following statement carries an
inhabitance conjunct.

\begin{theorem}[Setoid quotient classification]
\label{thm:quotient-contractible}
For every setoid \(S\) on \(X\),
\[
 \isContr(X/S)
 \quad\Longleftrightarrow\quad
 \bigl(\mathsf{Inh}(X)\ \wedge\ \SetTot(S)\bigr).
\tag{5.1}\label{eq:quotient-contractible}
\]
\end{theorem}

\begin{proof}
If \(X/S\) is contractible it is inhabited, and quotient induction on an
inhabitant exposes a representative, giving \(\mathsf{Inh}(X)\); no
representative is chosen for every class.  For \(x,y:X\), contractibility gives
\(\eta_S(x)=\eta_S(y)\), and quotient \emph{exactness} turns this back into
\(S(x,y)\).  Conversely, take \(\eta_S(x_0)\) as center for \(x_0\) supplied by
inhabitance; quotient induction reduces an arbitrary class to \(\eta_S(x)\), and
totality with quotient \emph{soundness} gives \(\eta_S(x_0)=\eta_S(x)\).
\end{proof}

The two directions use the two halves of quotient exactness, and this is the
only place either is needed.

\begin{corollary}[Chosen-center form]
\label{cor:quotient-total}
If \(x_0:X\), then \(\isContr(X/S)\Longleftrightarrow\SetTot(S)\).
\end{corollary}

\begin{warning}[The empty carrier]
\label{warn:empty-carrier}
For \(X=\varnothing\), \(\SetTot(S)\) holds vacuously while \(X/S\) is empty and
not contractible, so inhabitance cannot be dropped.  In the metadata application
the chosen \(m_0\) supplies it; it must not be silently inferred elsewhere.
\end{warning}

\subsection{Classification of metadata repairs}

Let
\[
 S_{b,h}\ \text{ be a setoid on }\ M(b,h)
\tag{5.2}\label{eq:setoid-family}
\]
for every \(b:A\) and \(h:a=b\).  Define the quotient metadata family and its
refined equality by
\[
 \overline M_S(b,h)=M(b,h)/S_{b,h},
 \qquad
 \overline R_S(a,b)=\sum_{h:a=b}\overline M_S(b,h),
\tag{5.3}\label{eq:quotient-metadata}
\]
with base
\[
       \overline r^S_a=(\refl_a,\eta_{S_{a,\refl_a}}(m_0)).
\tag{5.4}\label{eq:quotient-metadata-base}
\]

\begin{theorem}[Universal setoid-repair criterion]
\label{thm:universal-repair}
The quotient-repaired refinement admits unrestricted based elimination with
propositional beta exactly when the setoid on the reflexivity metadata fiber is
total:
\[
 \Jbeta_{\overline R_S}(a,\overline r^S_a)
 \quad\Longleftrightarrow\quad
 \SetTot(S_{a,\refl_a}).
\tag{5.5}\label{eq:universal-repair}
\]
\end{theorem}

\begin{proof}
Apply the metadata-fiber criterion to \(\overline M_S\):
\[
 \Jbeta_{\overline R_S}(a,\overline r^S_a)
 \quad\Longleftrightarrow\quad
 \isContr\!\left(M(a,\refl_a)/S_{a,\refl_a}\right).
\tag{5.6}\label{eq:repair-first-step}
\]
The raw reflexivity fiber is inhabited by \(m_0\).  Corollary
\ref{cor:quotient-total} identifies the right-hand side with
\(\SetTot(S_{a,\refl_a})\).
\end{proof}

The theorem is both a construction and an impossibility result.  The forward
direction says that every successful ordinary setoid repair must relate
\emph{every} pair \(x,y:M(a,\refl_a)\).  By quotient soundness,
\[
             \eta_{S_{a,\refl_a}}(x)
             =
             \eta_{S_{a,\refl_a}}(y)
             \qquad\text{for all }x,y.
\tag{5.7}\label{eq:no-visible-reflexivity}
\]
Therefore:
\begin{quote}
\centering\bfseries
No set-level quotient supporting unrestricted \(J^\beta\) can retain a
visible reflexivity-metadata distinction.
\end{quote}
This does not prohibit retaining raw metadata behind an opaque interface or
restricting motives.  It says that the quotient \emph{as the identity-like
object} has only one reflexivity class.

\begin{proposition}[Family-wide transport]
\label{prop:family-total}
If \(S_{a,\refl_a}\) is total, then every \(S_{b,h}\) is total.  Moreover, the
chosen \(m_0\) transports to an inhabitant of every \(M(b,h)\), so every
quotient \(M(b,h)/S_{b,h}\) is contractible.
\end{proposition}

\begin{proof}
Fix \(b:A\) and \(h:a=b\).  Equality induction on \(h\) reduces both assertions
to the reflexive case.  Totality is then the hypothesis, and inhabitance is
provided by \(m_0\).  Apply Corollary~\ref{cor:quotient-total} to the resulting
fiber.
\end{proof}

This proposition explains the precise sense in which the reflexivity test is
family-wide.  It is not a pointwise cardinality assertion about arbitrary
indices.  Equality induction transports the setoid family itself.  Empty
fibers cannot appear in this equality-indexed family once \(m_0\) is given,
whereas Warning~\ref{warn:empty-carrier} still governs standalone quotient
claims.

\subsection{A nontrivial quotient that repairs nothing}
\label{subsec:parity}

The necessity direction of Theorem~\ref{thm:universal-repair} has content only
if some genuine quotient fails it.  It is worth exhibiting one, because the
informal instruction ``quotient the administrative metadata'' can easily
produce a quotient that collapses a great deal and still repairs nothing.

\begin{example}[Trace-length parity]
\label{ex:parity}
Let \(M_{\mathrm{tr}}(b,h)=\List(\Step(A))\) be the trace metadata of
Section~\ref{sec:traces} and let \(S^{\mathrm{par}}\) relate two traces exactly
when their lengths have the same parity.  This setoid is far from discrete: each
class contains infinitely many pairwise distinct traces, and, for instance,
\[
 \bigl[\Empty\bigr]
 =
 \bigl[\,[\,a\xrightarrow{\refl_a}a,\;a\xrightarrow{\refl_a}a\,]\,\bigr].
\tag{5.8}\label{eq:parity-collapse}
\]
It is nevertheless not total, because \(\Empty\) has even length while
\([\,a\xrightarrow{\refl_a}a\,]\) has odd length.  These two traces therefore
have distinct images, so by Theorem~\ref{thm:quotient-contractible} the
reflexivity quotient is not contractible, and by
Theorem~\ref{thm:universal-repair} the parity-repaired refinement still admits
no unrestricted based eliminator.
\end{example}

Example~\ref{ex:parity} isolates the misconception cleanly.  Quotienting is not
measured by how much it identifies, nor by whether the resulting classes have
canonical representatives, nor by decidability of the induced equality.  The
only relevant question is whether the reflexivity quotient is contractible.
Parity identifies infinitely many traces and still fails; the indiscrete
relation of Section~\ref{subsec:indiscrete} is the smallest strengthening that
cannot fail, because by Theorem~\ref{thm:universal-repair} success already
forces every pair to be related.

\subsection{The indiscrete repair}
\label{subsec:indiscrete}

\begin{definition}[Indiscrete setoid]
\label{def:indiscrete}
The indiscrete setoid \(\nabla_X\) on \(X\) is
\[
                        \nabla_X(x,y)=\top
                        \qquad(x,y:X).
\tag{5.9}\label{eq:indiscrete}
\]
\(\top\) denotes the always-true proposition; it is not the data type \(\Unit\).
\end{definition}

The relation \(\nabla_X\) is total, hence its quotient is contractible whenever
\(X\) is inhabited.  Applied fiberwise, it is the canonical unconditional
repair supplied by the data \(m_0\).

\subsection{Exact factorization direction}

The indiscrete repair is canonical, and it is worth saying in which direction.
For setoids \(S,T\) on \(X\), call \(F:X/S\to X/T\) a
\emph{quotient-map-preserving factor} when \(F\circ\eta_S=\eta_T\).

\begin{theorem}[Relation inclusion versus quotient factorization]
\label{thm:factor-inclusion}
A quotient-map-preserving factor \(X/S\to X/T\) exists exactly when
\(S\subseteq T\), and it is then unique as a function.
\end{theorem}

\begin{proof}
Given \(S\subseteq T\), quotient recursion sets \(F(\eta_S(x))=\eta_T(x)\); the
inclusion is precisely the well-definedness obligation, by soundness for \(T\).
Conversely \(S(x,y)\) gives \(\eta_S(x)=\eta_S(y)\), hence
\(\eta_T(x)=\eta_T(y)\) after applying \(F\), and exactness returns \(T(x,y)\).
Uniqueness is quotient induction on a representative.
\end{proof}

Every \(S\) is contained in \(\nabla_X\), so arrows point \emph{toward} the
indiscrete quotient: it is terminal, not initial, in this
quotient-map-preserving category.  When \(S\) is itself total the inclusion
reverses as well, and quotient induction makes the two maps inverse,
\[
 X/S\Eqv X/\nabla_X.
\tag{5.10}\label{eq:total-indiscrete-equiv}
\]
So every successful setoid repair is, on the reflexivity fiber, canonically
equivalent to the indiscrete one.  This does not say the raw setoids are
definitionally equal, nor that discarded metadata can be recovered.

\begin{remark}[Maximal, minimal, terminal]
\label{rem:terminal-not-initial}
Three descriptions of \(\nabla_X\) are easily conflated.  It is
\emph{maximal} as a relation, \emph{minimal} in retained information, and
\emph{terminal} for quotient-map-preserving factors.  The first two are the same
fact seen from opposite sides; neither justifies an initiality claim.
\end{remark}

\begin{remark}[Scope of universality]
\label{rem:quotient-scope}
The necessity direction uses exactness,
\(\eta_S(x)=\eta_S(y)\Rightarrow S(x,y)\).  The classification therefore covers
ordinary exact setoid quotients in the present metatheory, and is not asserted
for truncations, higher quotients, opaque implementations without exactness, or
homotopical localization.
\end{remark}

\section{Projection kernels}
\label{sec:projection-kernel}

The metadata presentation \(\sum_hM(b,h)\) is convenient but not fundamental.
What matters is a relation carrying a reflexive projection to equality.

\subsection{Reflexive equality projections}

\begin{definition}[Reflexive equality projection]
\label{def:reflexive-projection}
Fix \(a:A\).  A reflexive equality projection consists of
\[
 R:A\to\U,\qquad r_0:R(a),\qquad
 q_b:R(b)\to(a=b),
\tag{6.1}\label{eq:reflexive-projection}
\]
together with \(q_a(r_0)=\refl_a\).
\end{definition}

\begin{definition}[Reflexivity kernel]
\label{def:reflexivity-kernel}
The reflexivity kernel of \(q\) is
\[
       \Ker(q)
       =
       \sum_{r:R(a)}\bigl(q_a(r)=\refl_a\bigr),
\tag{6.2}\label{eq:reflexivity-kernel}
\]
with center \((r_0,q_a(r_0)=\refl_a)\).
\end{definition}

\begin{theorem}[Projection/kernel equivalence]
\label{thm:projection-kernel-equiv}
There is a specified equivalence
\[
                       \sum_{b:A}R(b)\Eqv\Ker(q).
\tag{6.3}\label{eq:projection-kernel-equiv}
\]
\end{theorem}

\begin{proof}
Given \((b,r):\sum_bR(b)\), eliminate on \(q_b(r):a=b\).  In the reflexive case
return \(r:R(a)\), paired with proof irrelevance
\(q_a(r)=\refl_a\).  Conversely, forget the kernel equation and send
\((r,\rho)\) to \((a,r)\).  Equality induction on \(q_b(r)\) proves the first
inverse law; subtype extensionality and proof irrelevance prove the second.
\end{proof}

\begin{theorem}[Projection/kernel classification]
\label{thm:projection-kernel}
The relation row \(R\), based at \(r_0\), admits unrestricted based elimination
with propositional beta exactly when its reflexivity kernel is contractible:
\[
                  \Jbeta_R(a,r_0)
                  \quad\Longleftrightarrow\quad
                  \isContr(\Ker(q)).
\tag{6.4}\label{eq:projection-kernel}
\]
\end{theorem}

\begin{proof}
The total-space classification gives
\(\Jbeta_R(a,r_0)\leftrightarrow\isContr(\sum_bR(b))\).  Transport
contractibility across Theorem~\ref{thm:projection-kernel-equiv}.
\end{proof}

The metadata theorem is recovered by taking
\[
 R(b)=\sum_{h:a=b}M(b,h),
 \qquad
 q_b(h,m)=h.
\tag{6.5}\label{eq:metadata-as-projection}
\]
The kernel is then equivalent to \(M(a,\refl_a)\).  The projection theorem is
strictly more reusable: \(R\) need not be presented as a sigma type, and \(q\)
may forget an implementation-specific representation.

\section{The largest explicitly supported class of motives}
\label{sec:factorization}

Failure of unrestricted elimination does not mean that no useful elimination
principle remains.  The projection~\eqref{eq:projection} supports precisely the
class we can state and prove without inspecting metadata.

\begin{definition}[Strict equality-projection motive]
\label{def:strict-factor}
Given
\[
                  D:\prod_{b:A}(a=b)\longrightarrow\U,
\]
its pullback along \(\proj\) is the motive
\[
       C_D(b,h,m)=D(b,h)
       \qquad ((h,m):R_M(a,b)).
\tag{7.1}\label{eq:strict-factor}
\]
\end{definition}

\begin{theorem}[Motive factorization]
\label{thm:motive-factorization}
Every strict equality-projection motive \(C_D\) admits elimination
\[
  J_{\proj}:
  D(a,\refl_a)\longrightarrow
  \prod_{b:A}\prod_{p:R_M(a,b)}D(b,\proj p)
\tag{7.2}\label{eq:factored-j}
\]
with propositional beta at \(r^M_a\):
\[
                  J_{\proj}(d,a,r^M_a)=d.
\tag{7.3}\label{eq:factored-beta}
\]
More generally, the conclusion holds for any motive \(C\) equipped, for every
\(x=(b,h,m):T_M(a)\), with a specified fiberwise equivalence
\[
                         C(x)\Eqv D(b,h).
\tag{7.4}\label{eq:fiberwise-factor}
\]
\end{theorem}

\begin{proof}
For the strict case, ignore \(m\) and eliminate on \(h:a=b\):
\[
       J_{\proj}(d,b,(h,m))
       =\mathsf J_{=}(D,d,b,h):D(b,h).
\]
At \(h=\refl_a\), equality beta gives~\eqref{eq:factored-beta}.

For the fiberwise case, let \(e_x:C(x)\Eqv D(b,h)\) be the chosen equivalence.
Map the base datum \(c:C(c_M)\) forward along \(e_{c_M}\), apply the strict
construction in \(D\), and map the result back along \(e_x^{-1}\).  At the
center the equality eliminator has propositional beta, and
\(e_{c_M}^{-1}(e_{c_M}(c))=c\), proving the required beta equation.
\end{proof}

The theorem identifies the largest class \emph{explicitly supported here}:
motives supplied with a factorization through equality, strictly or up to
fiberwise equivalence.  We do not claim that this is a semantic maximality
theorem about isolated motives.  A particular metadata-sensitive motive might
have a section for accidental reasons.  What fails is one uniform operator
covering every motive from one base datum.

\begin{corollary}[Metadata irrelevance for factored elimination]
\label{cor:factored-irrelevance}
For \(p=(h,m)\) and \(q=(k,n)\) in \(R_M(a,b)\), if \(h=k\), then factored
elimination gives propositionally equal results after transport along that
equality.  In the proof-irrelevant ambient setting this applies to every
pair \(p,q:R_M(a,b)\).
\end{corollary}

\subsection{Two remedies}

\paragraph{Restrict motives.}
Keep the full record \(R_M\), but expose only motives of the form
\eqref{eq:strict-factor} or~\eqref{eq:fiberwise-factor}.  This retains
histories, costs, and traces for inspection by nondependent operations while
preventing a purported identity eliminator from distinguishing them.

\paragraph{Hide or quotient metadata.}
Replace \(M\) by a family \(\overline M\) whose reflexivity fiber is
contractible.  Examples include an opaque interface with a single observable
value, a propositional truncation when inhabited, or an explicit quotient
collapsing all reflexive histories.  Theorem~\ref{thm:metadata-fiber} then
provides unrestricted elimination.  This is not a proof that the original
metadata was irrelevant; it is a new identity-like object in which the
distinction has been removed.

\begin{warning}[Quotienting must match the intended observation]
Merely quotienting by administrative rewrites is insufficient if two
reflexive metadata classes remain.  The criterion asks for contractibility of
the resulting reflexivity fiber, not merely confluence, normalization, or
decidable equality.
\end{warning}

\section{Beyond lists: a classification guide}
\label{sec:examples}

The criterion applies to any observable decoration.  The following examples
are not encodings of one another; they illustrate which property matters.

\begin{table}[t]
\centering
\small
\renewcommand{\arraystretch}{1.22}
\begin{tabular}{>{\raggedright\arraybackslash}p{0.20\textwidth}
                >{\raggedright\arraybackslash}p{0.36\textwidth}
                >{\raggedright\arraybackslash}p{0.34\textwidth}}
\toprule
Metadata & Reflexivity fiber & Unrestricted \(J^\beta\)\\
\midrule
Hidden token
 & A singleton or inhabited subsingleton
 & Yes: the fiber is contractible.\\
Derivation history
 & Empty history and a visible reflexive inference
 & No, unless histories are hidden or collapsed.\\
Cost
 & Distinct costs \(0\) and \(1\) for reflexive computations
 & No; uniqueness of a minimum is irrelevant.\\
Certificate
 & An inhabited proof-irrelevant proposition
 & Yes in the present metatheory.\\
Certificate as data
 & Two distinguishable certificate objects
 & No, even if both certify the same fact.\\
Normalization log
 & Empty log and a nonempty administrative loop
 & No; unique normal forms do not imply unique logs.\\
Proof label
 & Exactly one reflexivity label
 & Yes; two visible reflexivity labels make it fail.\\
\bottomrule
\end{tabular}
\caption{The reflexivity fiber, rather than the informal purpose of metadata,
determines unrestricted based elimination.}
\label{tab:examples}
\end{table}

Table~\ref{tab:examples} is read off the reflexivity fiber and nothing else,
and three of its rows are worth a comment because they are commonly misread.

\emph{Costs} are not rescued by structure: additivity under composition,
decidable comparison, and uniqueness of an optimum are all irrelevant, because
optimization supplies an order and a preferred element, not an equality between
all costs.  Exposing instead the \emph{proposition} that the stored cost is
optimal does restore elimination, an inhabited proposition being contractible
here.

\emph{Certificates} hide a representation choice.  An inhabited proof-irrelevant
certificate fiber is contractible, so elimination holds
(Corollary~\ref{cor:contractible-decoration}); certificates carried as records
with serial numbers, proof trees, or signatures may certify the same equality
while remaining distinguishable, and then it fails
(Corollary~\ref{cor:visible-alternatives}).

\emph{Normalization logs} separate three properties that are easily conflated:
uniqueness of normal forms, connectedness of logs under log rewrites, and
equality of all observable reflexive logs.  Only the third is the metadata-fiber
condition; the second implies it only after quotienting, at which point
Theorem~\ref{thm:universal-repair} applies to the quotient.  The same pattern
governs derivation histories, where an empty derivation and an explicit
reflexivity-rule derivation obstruct elimination while they remain visible as
syntax, and proof labels, where arbitrarily elaborate indexing away from
reflexivity is harmless because the total-space equivalence transports
everything into one fiber.

\part{The computational-path rewrite quotient}
\label{part:case-study}

Part~\ref{part:case-study} applies the classification to one concrete
implementation and finds that it collapses.  The general theorems of
Part~\ref{part:general} are not tied to this implementation; the strongest
concrete result very much is.


\section{Computational traces: the principal case study}
\label{sec:traces}

\subsection{The trace-carrying record}

For a carrier \(A\), let an elementary step be
\[
       \Step(A)=\sum_{x,y:A}(x=y).
\tag{9.1}\label{eq:step}
\]
A step stores source, target, and an ambient equality witness.  Define
\[
       \PathTy_A(a,b)
       =\List(\Step(A))\times(a=b),
\tag{9.2}\label{eq:path}
\]
written as a record \((L,h)\).  Its reflexive element is
\[
                         \mathbf r_a=(\Empty,\refl_a).
\tag{9.3}\label{eq:empty-refl}
\]
The list is observable data.  Concatenation and reversal may be used to define
composition and symmetry, but those operations are not needed for the
classification.

Set
\[
 M_{\mathrm{tr}}(b,h)=\List(\Step(A)).
\tag{9.4}\label{eq:trace-metadata}
\]
Reordering fields gives mutually inverse maps
\[
\begin{aligned}
 S:\PathTy_A(a,b)&\longrightarrow
      \sum_{h:a=b}M_{\mathrm{tr}}(b,h),
 & S(L,h)&=(h,L),\\
 T:\sum_{h:a=b}M_{\mathrm{tr}}(b,h)&\longrightarrow
      \PathTy_A(a,b),
 & T(h,L)&=(L,h).
\end{aligned}
\tag{9.5}\label{eq:path-refinement}
\]
Both inverse equations are reflexivity of records.  Thus this is not only an
analogy: \(\PathTy\) is the corresponding equality refinement.

The reflexivity fiber can be computed outright, which makes the obstruction
below quantitative rather than merely nonempty.

\begin{proposition}[The trace fiber, computed]
\label{prop:trace-fiber-computed}
A step is determined by its source, so \(\Step(A)\Eqv A\).  Consequently
\[
 \PathTy_A(a,b)\Eqv\List(A)\times(a=b),
 \qquad
 \PathTy_A(a,a)\Eqv\List(A).
\tag{9.6}\label{eq:trace-fiber-computed}
\]
\end{proposition}

\begin{proof}
Send \(s:\Step(A)\) to its source, and \(x:A\) to \((x,x,\refl_x)\).
Eliminating the stored equality of \(s\) replaces the target by the source, and
the equality field is then proof-irrelevant, so the two composites are
identities.  Applying this pointwise to lists gives
\(\List(\Step(A))\Eqv\List(A)\), and reordering fields as in
\eqref{eq:path-refinement} gives the first equivalence.  For the second, the
remaining component \(a=a\) is a proof-irrelevant proposition, so two loops with
the same trace are equal.
\end{proof}

\begin{definition}[Two reflexive traces]
\label{def:two-traces}
Besides \(\mathbf r_a\), write
\(\mathbf s_a=\bigl([\,(a,a,\refl_a)\,],\refl_a\bigr):\PathTy_A(a,a)\) for the
singleton reflexive trace.  List no-confusion gives
\(\mathbf r_a\ne\mathbf s_a\).
\end{definition}

\begin{theorem}[No unrestricted trace-sensitive \(J\)]
\label{thm:no-trace-j}
For every \(A\) and \(a:A\), the family \(\PathTy_A(a,-)\), based at
\(\mathbf r_a\), has no unrestricted based eliminator satisfying propositional
beta.
\end{theorem}

\begin{proof}
By~\eqref{eq:path-refinement} the record is the equality refinement with
metadata~\eqref{eq:trace-metadata}, whose reflexivity fiber
\(\List(\Step(A))\) contains \(\mathbf r_a\ne\mathbf s_a\) and so is not
contractible; apply Theorem~\ref{thm:metadata-fiber}.
\end{proof}

Proposition~\ref{prop:trace-fiber-computed} makes the failure
\emph{unconditional}: no hypothesis on \(A\) can rescue the raw record, because
a pointed carrier already forces two visible reflexive traces.  This is the
contrast with the \(\PRW\) criterion of Theorem~\ref{thm:raw-criterion}, which
is conditional and which Section~\ref{sec:collapse} then discharges in the
opposite direction.

One may recover the familiar direct counter-motive by taking \(C(b,L,h)=\Unit\)
when \(L=\Empty\) and \(\varnothing\) otherwise: it has a base value but none at
\(\mathbf s_a\).  That is a useful diagnostic, not the classification proof.

\subsection{Composability is not the issue}

The raw list in~\eqref{eq:path} need not enforce adjacency.  It is tempting to
attribute Theorem~\ref{thm:no-trace-j} to that permissiveness.  The following
intrinsic version rules out the diagnosis.

\begin{definition}[Composable visible trace]
\label{def:composable}
A list \([s_1,\ldots,s_n]\) is composable from \(a\) to \(b\) when the source
of \(s_1\) is \(a\), the target of \(s_i\) is the source of \(s_{i+1}\), and
the target of \(s_n\) is \(b\).  The empty list is composable from \(a\) to
\(a\).  Write \(\mathsf{CTr}_A(a,b)\) for the subtype of lists satisfying this
condition.
\end{definition}

Both
\[
      \epsilon_a=(\Empty,\text{empty-chain proof})
      \quad\text{and}\quad
      \sigma_a=([\,a\xrightarrow{\refl_a}a\,],
                   \text{singleton-chain proof})
\tag{9.7}\label{eq:composable-two}
\]
belong to \(\mathsf{CTr}_A(a,a)\).  Equality of subtype values would imply
equality of their underlying lists, so \(\epsilon_a\ne\sigma_a\).

\begin{corollary}[Composable-trace obstruction]
\label{cor:composable}
The equality refinement
\[
   \sum_{h:a=b}\mathsf{CTr}_A(a,b),
\tag{9.8}\label{eq:composable-refinement}
\]
based at \((\refl_a,\epsilon_a)\), has no unrestricted based eliminator with
propositional beta.
\end{corollary}

\begin{proof}
The reflexivity metadata fiber
\(\mathsf{CTr}_A(a,a)\) contains the distinct values
\(\epsilon_a\) and \(\sigma_a\).  It is noncontractible, so
Theorem~\ref{thm:metadata-fiber} applies.
\end{proof}

The singleton step starts at \(a\), ends at \(a\), and has reflexive composite.
It survives any condition requiring adjacency, correct outer endpoints, and a
composite equality matching \(\refl_a\).  The obstruction is therefore:
\[
 \boxed{\text{visible noncanonical reflexive metadata, not
 non-composability.}}
\tag{9.9}\label{eq:boxed-obstruction}
\]

\subsection{What remains available}

The endpoint-equality projection \((L,h)\mapsto h\) is
exactly~\eqref{eq:projection}, so for \(D:\prod_{b:A}(a=b)\to\U\)
Theorem~\ref{thm:motive-factorization} supplies
\[
 D(a,\refl_a)\longrightarrow
 \prod_{b:A}\prod_{P:\PathTy_A(a,b)}D(b,\proj P)
\tag{9.10}\label{eq:trace-factored-j}
\]
with propositional beta at \(\mathbf r_a\), independently of the list.  Trace
algebra---concatenation, inversion, rewrite equivalence---therefore coexists
with equality-component elimination.  What cannot coexist is an unrestricted
eliminator whose motive may separate \(\mathbf r_a\) from \(\mathbf s_a\).

\subsection{Rewrite evidence, and the quotient}
\label{subsec:raw-versus-rweq}

Composition of records concatenates their lists, and \(\mathsf{RwEq}\) records
the LND\(_{\mathrm{EQ}}\)-TRS coherences; for a two-step path \(p\) one has the
certificate \(p\cdot p^{-1}\mathrel{\mathsf{RwEq}}\mathbf r\).  Such a
certificate organizes traces.  It does not, by itself, equate every list as
record data, and in particular it does not make the raw reflexivity fiber
contractible.

The two witnesses of Theorem~\ref{thm:no-trace-j} illustrate the gap.  Writing
\(\epsilon_a=\mathbf r_a\) for the empty reflexive record and
\(\sigma_a=\mathbf s_a\) for the singleton one, their step lists have lengths
zero and one, so \(\epsilon_a\ne\sigma_a\) as records; yet the primitive
transport-beta rewrite at a constant family gives
\[
 \sigma_a\mathrel{\mathsf{RwEq}}\epsilon_a,
 \qquad\text{hence}\qquad
 [\epsilon_a]=[\sigma_a]\ \text{in }\PRW_A(a,a).
\tag{9.11}\label{eq:raw-empty-singleton-rweq}
\]
A motive on the raw record may inspect the list and separate the two; a motive on
\(\PRW\) sees only the class, and this particular distinction has been removed by
an actual derivation.  Equality of the outputs of the implemented normalization
function would not establish~\eqref{eq:raw-empty-singleton-rweq}---that function
depends only on \(\pi_{=}\)---so the argument uses the named primitive rule
instead.  Whether \emph{all} distinctions vanish is the local-K question of
Theorem~\ref{thm:pathrw-k}, settled affirmatively, for a quite different reason,
in Section~\ref{sec:collapse}.

\section{\texorpdfstring{The \(\mathsf{PathRwQuot}\) criterion}{The PathRwQuot criterion}}
\label{sec:pathrwquot-criterion}

\subsection{Applying the kernel theorem}

Let
\[
 R_{\mathsf{rw}}(b)=\PRW_A(a,b)
 =\PathTy_A(a,b)/\mathsf{RwEq},
\tag{10.1}\label{eq:pathrw-row}
\]
based at the class \([\mathbf r_a]\).  The quotient operation
\[
 \mathsf{toEq}_b:\PRW_A(a,b)\longrightarrow(a=b)
\tag{10.2}\label{eq:pathrw-toeq}
\]
forgets rewrite traces and is well-defined because every primitive rewrite
preserves the underlying equality proof.

The reflexivity kernel is
\[
 \Ker(\mathsf{toEq})
 =
 \sum_{x:\PRW_A(a,a)}
       \bigl(\mathsf{toEq}(x)=\refl_a\bigr).
\tag{10.3}\label{eq:pathrw-kernel}
\]
Ambient equality proofs live in \(\mathsf{Prop}\), so the second component is
automatic by proof irrelevance.  Projection to the first component and
reinsertion of that proof give
\[
              \Ker(\mathsf{toEq})\Eqv\PRW_A(a,a).
\tag{10.4}\label{eq:pathrw-kernel-loop}
\]

\begin{definition}[Local quotient-level axiom K]
\label{def:local-k}
Define
\[
 \LocalK(A,a)
 \quad:\!\!\Longleftrightarrow\quad
 \prod_{x:\PRW_A(a,a)}x=[\mathbf r_a].
\tag{10.5}\label{eq:local-k}
\]
\end{definition}

This is a statement about classes of computational paths.  It is not the
assertion that two raw records with different step lists are equal, and
it is not obtained merely by observing that their \(\mathsf{toEq}\) values are
proof-irrelevant.

\begin{theorem}[\(\mathsf{PathRwQuot}\)/K classification]
\label{thm:pathrw-k}
For every \(A\) and \(a:A\), the following are equivalent:
\[
\begin{split}
 &\text{\(\PRW_A(a,-)\), based at \([\mathbf r_a]\), admits
 unrestricted \(J^\beta\);}\\
 &\isContr\bigl(\PRW_A(a,a)\bigr);\\
 &\LocalK(A,a).
\end{split}
\tag{10.6}\label{eq:pathrw-k}
\]
\end{theorem}

\begin{proof}
The first equivalence is Theorem~\ref{thm:projection-kernel} followed by
\eqref{eq:pathrw-kernel-loop}.  For the second, a contraction identifies every
loop class with the reflexive class.  Conversely, the equations in
\(\LocalK(A,a)\), reversed in orientation, contract the loop quotient at
\([\mathbf r_a]\).
\end{proof}

\begin{corollary}[Global form]
\label{cor:pathrw-global-k}
The family \(\PRW_A(a,-)\) admits such an eliminator for every \(a:A\) exactly
when \(\prod_{a:A}\LocalK(A,a)\).
\end{corollary}

\begin{remark}[The criterion is exact, in both directions]
\label{rem:rweq-no-slogan}
Theorem~\ref{thm:pathrw-k} says neither that quotienting by \(\mathsf{RwEq}\)
always repairs unrestricted elimination nor that it always fails; it is the
exact criterion, and a concrete rule set must be checked against it.
Section~\ref{sec:collapse} performs that check for the formalized rules.
\end{remark}

\begin{remark}[Normalization is only one direction]
\label{rem:rweq-normalization}
\(\mathsf{RwEq}\) is the equivalence closure generated by explicit primitive
rewrites.  The implemented normalization function, which returns the one-step
chain of a path's underlying equality, satisfies
\(p\mathrel{\mathsf{RwEq}}q\Rightarrow\mathsf{normalize}(p)=\mathsf{normalize}(q)\),
but the converse is not part of the definition and cannot replace an
\(\mathsf{RwEq}\) witness; a normal-form characterization needs a separately
supplied complete step system.  Every derivation below therefore exhibits named
primitive rules.
\end{remark}

\subsection{Removing the quotient from the criterion}
\label{subsec:raw-criterion}

Theorem~\ref{thm:pathrw-k} is stated at quotient level, which makes it awkward
to check: verifying \(\LocalK(A,a)\) appears to require reasoning about
classes.  It can be pushed back to raw paths.  Doing so also exhibits
Section~\ref{sec:projection-kernel} as an \emph{instance} of
Section~\ref{sec:setoid-repair} rather than a parallel development, because
\(\PRW_A(a,b)\) is literally an ordinary setoid quotient: its carrier is
\(\PathTy_A(a,b)\) and its relation is \(\mathsf{RwEq}\).

\begin{definition}[Rewrite-connected loops]
\label{def:rweq-total}
Say that \(\mathsf{RwEq}\) is \emph{total on loops at \(a\)} when every two raw
loops admit a rewrite certificate:
\[
 \RwTot(A,a)
 \quad:\!\!\Longleftrightarrow\quad
 \prod_{p,q:\PathTy_A(a,a)}
   \bigl\|\,p\mathrel{\mathsf{RwEq}}q\,\bigr\|.
\tag{10.7}\label{eq:rweq-total}
\]
Here \(\|-\|\) is the propositional inhabitance used by the defining setoid, so
\eqref{eq:rweq-total} is exactly \(\SetTot\) for the relation \(\mathsf{RwEq}\)
on the carrier \(\PathTy_A(a,a)\).  No quotient occurs
in~\eqref{eq:rweq-total}.
\end{definition}

\begin{theorem}[Raw-level criterion]
\label{thm:raw-criterion}
For every \(A\) and \(a:A\), the following are equivalent:
\[
\begin{split}
 &\text{\(\PRW_A(a,-)\), based at \([\mathbf r_a]\), admits
   unrestricted \(J^\beta\);}\\
 &\isContr\bigl(\PRW_A(a,a)\bigr);\qquad
  \LocalK(A,a);\qquad
  \RwTot(A,a).
\end{split}
\tag{10.8}\label{eq:raw-criterion}
\]
\end{theorem}

\begin{proof}
Equivalence of the first three is Theorem~\ref{thm:pathrw-k}.  For the last,
apply Corollary~\ref{cor:quotient-total} to the setoid \(\mathsf{RwEq}\) on the
carrier \(\PathTy_A(a,a)\); the chosen center \(\mathbf r_a\) supplies
inhabitance.  Sufficiency is quotient soundness and necessity is quotient
exactness, exactly as in Theorem~\ref{thm:quotient-contractible}.
\end{proof}

\begin{corollary}[Eliminators manufacture certificates]
\label{cor:eliminator-certificates}
If \(\PRW_A(a,-)\) admits unrestricted based elimination with propositional
beta, then for every pair of raw loops \(p,q:\PathTy_A(a,a)\) there is an
actual \(\mathsf{RwEq}\) derivation between them.
\end{corollary}

Corollary~\ref{cor:eliminator-certificates} is the sharp form of the ``no
unconditional slogan'' point.  Assuming such an eliminator is not a mild
structural convenience: it is exactly as strong as asserting that the rewrite
system connects every raw loop at \(a\) to every other.  Conversely, this is
what one proves in practice.  Proposition~\ref{prop:subsingleton-loop-quotient}
below establishes \(\RwTot(A,a)\) directly, by constructing the derivations, and
only then reads off contractibility; it does not first form the quotient.

A small instance shows the shape of such a derivation.  Write
\(\mathbf s_a=\bigl([\,a\xrightarrow{\refl_a}a\,],\refl_a\bigr)\) for the
singleton reflexive trace, formally introduced in
Definition~\ref{def:two-traces}.  Composing it with itself gives a two-step raw
loop, and
\[
 \mathbf s_a\cdot\mathbf s_a
 \;\mathrel{\mathsf{RwEq}}\;
 \mathbf r_a
\tag{10.9}\label{eq:double-singleton}
\]
is obtained by congruence under composition, the primitive
transport-reflexivity erasure, the left-unit rule, and a second erasure.  The
two sides of~\eqref{eq:double-singleton} are distinct \(\PathTy\) records, so
this is a genuine derivation and not an instance of reflexivity.

\begin{remark}[Global form]
\label{rem:raw-global}
Quantifying over base points, family-wide elimination for \(\PRW\) is
equivalent to \(\prod_{a:A}\RwTot(A,a)\).  Theorem~\ref{thm:raw-criterion}
makes the remaining question completely elementary: it asks only whether two
given step lists are connected by primitive rewrites.  In an earlier version of
this work we left that question open, since it concerns the specific primitive
rule set of LND\(_{\mathrm{EQ}}\)-TRS rather than the classification.
Section~\ref{sec:collapse} settles it, and the answer is affirmative for every
carrier.
\end{remark}

The contrast with the unquotiented record is worth stating explicitly.  For raw
\(\PathTy\) the answer is unconditional and negative for every pointed carrier
(Theorem~\ref{thm:no-trace-j}, sharpened by
Proposition~\ref{prop:trace-fiber-computed}).  For \(\PRW\) the criterion is
\eqref{eq:rweq-total}, and Section~\ref{sec:collapse} verifies it outright.
Quotienting therefore changes the answer from ``always fails'' to ``always
succeeds'', and the classification tells us exactly what is paid for the
change.

\section{The collapse}
\label{sec:collapse}

Theorem~\ref{thm:raw-criterion} reduces a question about eliminators to a
question about the primitive rewrite rules.  This section answers it.  One rule
decides the matter, and it does so on every carrier at once.

\subsection{A trace-erasing packaging rule}

Two primitives of the artifact interact.  Congruence transports a trace
pointwise,
\[
 \mathsf{congrArg}(f,(L,h))
 =\bigl(\,\mathsf{map}\,(f_*)\,L,\ \mathrm{ap}_f(h)\,\bigr),
\tag{11.1}\label{eq:congrarg-trace}
\]
whereas the packaging operation for function types records nothing at all:
\[
 \mathsf{lamCongr}
   \Bigl(\,\lambda x.\,P_x:\prod_{x:\alpha}\PathTy(f\,x,g\,x)\Bigr)
 =\bigl(\,[\,]\ ,\ \mathsf{funext}(\lambda x.\,\pi_{=}P_x)\,\bigr).
\tag{11.2}\label{eq:lamcongr-trace}
\]
Equation~\eqref{eq:lamcongr-trace} is a design decision, not a derived fact:
there is no canonical way to concatenate an \(\alpha\)-indexed family of traces
into a single trace, and the implementation therefore keeps none.

The primitive application rule of the rewriting system relates the two:
\[
 \mathsf{congrArg}\bigl(\lambda h.\,h\,x,\ \mathsf{lamCongr}(P)\bigr)
 \;\triangleright\;
 P_x .
\tag{11.3}\label{eq:fun-app-beta}
\]

\begin{lemma}[The collapsing rule]
\label{lem:collapse-rule}
Write \(\varepsilon_{a,b}=([\,],h)\) for the trace-free path determined by an
ambient equality \(h:a=b\); by proof irrelevance it does not depend on the
choice of \(h\).  Then for every carrier \(A\), all \(a,b:A\), and \emph{every}
\(p:\PathTy_A(a,b)\) there is a single primitive rewrite
\[
 \varepsilon_{a,b}\;\triangleright\;p .
\tag{11.4}\label{eq:collapse-step}
\]
\end{lemma}

\begin{proof}
Instantiate~\eqref{eq:fun-app-beta} at the unit domain \(\alpha=\mathbf 1\),
with \(f=\lambda\_.a\), \(g=\lambda\_.b\), \(P=\lambda\_.p\) and \(x=\star\).
By~\eqref{eq:lamcongr-trace} the inner packaging has empty trace, and
by~\eqref{eq:congrarg-trace} congruence maps the empty trace to the empty trace,
so the left-hand side of~\eqref{eq:fun-app-beta} is literally
\(([\,],\mathrm{ap}_{(\lambda h.h\star)}(\mathsf{funext}(\lambda\_.\pi_{=}p)))\).
Its two fields are the empty list and an ambient proof of \(a=b\), so it is
\(\varepsilon_{a,b}\).  The right-hand side is \(P_\star=p\).
\end{proof}

Lemma~\ref{lem:collapse-rule} is the whole story.  The trace-free path is a
common one-step predecessor of every path with the same endpoints, so the
rewrite relation cannot separate any two of them.

\subsection{Totality, and unrestricted elimination}

\begin{theorem}[\(\mathsf{RwEq}\) is total]
\label{thm:rweq-total}
For every carrier \(A\) and all \(a,b:A\), any two computational paths
\(p,q:\PathTy_A(a,b)\) are related:
\[
 p\;\mathrel{\mathsf{RwEq}}\;q .
\tag{11.5}\label{eq:rweq-total-thm}
\]
In particular \(\RwTot(A,a)\) holds for every pointed carrier.
\end{theorem}

\begin{proof}
By Lemma~\ref{lem:collapse-rule} there are single rewrites
\(\varepsilon_{a,b}\triangleright p\) and
\(\varepsilon_{a,b}\triangleright q\).  Compose the inverse of the first with
the second inside the equivalence closure.
\end{proof}

The derivation in Theorem~\ref{thm:rweq-total} is genuine: it has two stages and
passes through a third path, \(\varepsilon_{a,b}\), which is in general distinct
from both endpoints as a record.

\begin{corollary}[Unconditional repair]
\label{cor:unconditional-repair}
For every carrier \(A\) and every \(a,b:A\):
\begin{enumerate}[label=(\roman*),itemsep=1pt,topsep=3pt]
\item \(\PRW_A(a,b)\) is a subsingleton, and
  \(\PRW_A(a,b)\Eqv(a=b)\);
\item \(\isContr\bigl(\PRW_A(a,a)\bigr)\);
\item \(\LocalK(A,a)\) holds, and so does the global form;
\item \(\PRW_A(a,-)\), based at \([\mathbf r_a]\), admits unrestricted based
  elimination into every motive, with propositional beta.
\end{enumerate}
\end{corollary}

\begin{proof}
Theorem~\ref{thm:rweq-total} gives \(\SetTot\) for the defining setoid, so
quotient soundness identifies all classes, which is~(i); the map to \(a=b\) is
\(\pi_{=}\) and its inverse sends \(h\) to \([\varepsilon_{a,b}]\).  Items
(ii)--(iv) are then Theorem~\ref{thm:raw-criterion} read from right to left.
\end{proof}

So the answer to the question that motivates this paper is affirmative, in the
following precise sense.

\begin{theorem}[Path induction for computational paths]
\label{thm:quotient-path-induction}
Let \(A\) be any type and \(a:A\).  For every motive
\[
        C:\prod_{b:A}\PRW_A(a,b)\longrightarrow\U
\]
---including motives that inspect the quotient class itself---and every
\(d:C(a,[\mathbf r_a])\) there is
\[
 \mathcal J_{\PRW}(C,d):\prod_{b:A}\prod_{x:\PRW_A(a,b)}C(b,x),
 \qquad
 \mathcal J_{\PRW}(C,d,a,[\mathbf r_a])=d .
\tag{11.6}\label{eq:prw-J}
\]
The Martin-L\"of (unbased) form and the associated transport operator follow.
\end{theorem}

\begin{proof}
Given \(x:\PRW_A(a,b)\), the projection \(\pi_{=}x:a=b\) allows equality
elimination to reduce to \(b=a\); Corollary~\ref{cor:unconditional-repair}(iii)
then supplies \(x=[\mathbf r_a]\), along which \(d\) transports.
\end{proof}

\begin{remark}[The beta law is judgmental here]
\label{rem:judgmental-beta}
Everywhere else in this paper beta laws are propositional, and the
classification is stated for that weaker notion.  For the eliminator of
Theorem~\ref{thm:quotient-path-induction} the equation in~\eqref{eq:prw-J} in
fact holds by reflexivity in the formalization.  The reason is specific and does
not generalize: the transport is performed along an \emph{ambient equality
proof}, and definitional proof irrelevance identifies that proof with
reflexivity, after which the recursor reduces.  We flag this as a property of
the construction rather than as a strengthening of the classification theorems.
\end{remark}

\subsection{What has been bought, and what has been paid}

Corollary~\ref{cor:unconditional-repair} is the trade-off predicted by
Theorem~\ref{thm:universal-repair}, instantiated at the artifact's own rewrite
relation, and it is met in the extreme form: by~(i) the quotient retains the
ambient equality and nothing else, so \(\PRW\) is equivalent to the indiscrete
repair of Section~\ref{subsec:indiscrete}.  Two readings follow, and they should
be kept apart.

\begin{enumerate}[label=(\alph*)]
\item \emph{Positive.}  Raw records fail unrestricted \(J\) on every pointed
  carrier, but the rewrite quotient satisfies it on every carrier, with no
  hypothesis and no additional axiom.  A user of \(\PRW\) may reason by path
  induction---because \(\PRW\) is ambient equality.
\item \emph{Cautionary.}  For the same reason \(\PRW\) carries no rewrite
  information, and no invariant defined on it can be nonconstant.  Any
  construction whose interest depends on separating two classes of \(\PRW\) is,
  under the present rule set, vacuous.
\end{enumerate}

Section~\ref{sec:design} shows this is not a coincidence of these rules but a
theorem about ordinary quotients.  The dichotomy with the raw record is sharp.

\begin{theorem}[Raw versus quotient]
\label{thm:raw-vs-quotient}
For every pointed carrier \((A,a)\),
\(\neg\isContr(\PathTy_A(a,a))\) and \(\isContr(\PRW_A(a,a))\).
\end{theorem}

\begin{proof}
Proposition~\ref{prop:trace-fiber-computed} and
Corollary~\ref{cor:unconditional-repair}(ii).
\end{proof}

\subsection{Sharpness: the groupoid fragment does not collapse}
\label{subsec:sharpness}

It matters whether the collapse is intrinsic to computational-path rewriting or
an artifact of particular rules.  It is the latter, and the separating invariant
is elementary.

\begin{definition}[Groupoid fragment]
\label{def:groupoid-fragment}
Let \(\triangleright_{\mathrm{grp}}\) be the sub-relation of \(\triangleright\)
generated by the path-algebra rules only:
\(\mathsf{symm}\,\mathbf r\triangleright\mathbf r\),
double symmetry, the two unit laws, the two inverse laws, contravariance of
symmetry over composition, associativity, the two cancellation rules, and the
three structural congruences for \(\mathsf{symm}\) and composition.  Write
\(\mathrel{\mathsf{RwEq}}_{\mathrm{grp}}\) for its equivalence closure.
\end{definition}

\begin{proposition}[Parity is a groupoid invariant]
\label{prop:parity-invariant}
Let \(\ell(p)\) denote the length of the recorded trace of \(p\).  If
\(p\mathrel{\mathsf{RwEq}}_{\mathrm{grp}}q\) then
\(\ell(p)\equiv\ell(q)\pmod 2\).
\end{proposition}

\begin{proof}
Composition concatenates traces and symmetry reverses them, so
\(\ell(p\cdot q)=\ell(p)+\ell(q)\) and \(\ell(p^{-1})=\ell(p)\).  Each generator
is then checked directly.  The unit, associativity and contravariance rules
preserve \(\ell\) exactly.  The inverse laws send \(2\ell(p)\) to \(0\) and the
cancellation rules send \(2\ell(p)+\ell(q)\) to \(\ell(q)\); both preserve
parity.  The congruences preserve parity by induction.  Parity is then
preserved by reflexivity, symmetry and transitivity of the closure.
\end{proof}

\begin{theorem}[The collapse is caused by the function-space rules]
\label{thm:sharpness}
The empty and singleton reflexive traces satisfy
\[
 \mathbf r_a
 \mathrel{\mathsf{RwEq}}
 \mathbf s_a,
 \qquad\text{but}\qquad
 \neg\bigl(\mathbf r_a
   \mathrel{\mathsf{RwEq}}_{\mathrm{grp}}\mathbf s_a\bigr).
\tag{11.7}\label{eq:sharpness}
\]
Hence \(\triangleright_{\mathrm{grp}}\) is not total, its quotient is not
contractible, and it does not by itself repair unrestricted \(J\).
\end{theorem}

\begin{proof}
The first is Theorem~\ref{thm:rweq-total}.  For the second,
\(\ell(\mathbf r_a)=0\) and \(\ell(\mathbf s_a)=1\), so
Proposition~\ref{prop:parity-invariant} forbids a groupoid derivation.  The
remaining assertions follow from Theorem~\ref{thm:quotient-contractible} applied
to the setoid generated by \(\triangleright_{\mathrm{grp}}\), whose totality has
just been refuted.
\end{proof}

Theorem~\ref{thm:sharpness} also shows that the criterion of
Section~\ref{sec:projection-kernel} is doing work in both directions.  Writing
\(\PathTy_A(a,a)/\mathord{\mathsf{RwEq}}_{\mathrm{grp}}\) for the fragment's loop
quotient,
\[
 \neg\isContr\bigl(\PathTy_A(a,a)/\mathord{\mathsf{RwEq}}_{\mathrm{grp}}\bigr)
 \qquad\text{and}\qquad
 \isContr\bigl(\PRW_A(a,a)\bigr)
\tag{11.8}\label{eq:fragment-vs-full}
\]
for every pointed carrier.  Two rewrite systems over the same carrier, differing
only in the function-space rules, land on opposite sides of the classification.

Theorem~\ref{thm:sharpness} thus localises the phenomenon.  The classical
LND\(_{\mathrm{EQ}}\)-TRS coherences---the ones making paths a groupoid up to
rewriting---are compatible with a nontrivial rewrite quotient; what collapses it
is the combination of the function-space rules \eqref{eq:fun-app-beta} with the
trace-erasing convention \eqref{eq:lamcongr-trace}.
Section~\ref{sec:design} turns this observation into a constraint any
replacement rule set must satisfy.

\begin{remark}[Scope of the collapse theorem]
\label{rem:collapse-scope}
Theorem~\ref{thm:rweq-total} is a statement about the rule set formalized in the
accompanying artifact, and about the specific trace semantics of the operations
appearing in~\eqref{eq:lamcongr-trace}.  It is not a claim about
computational-path calculi in general, and in particular it does not apply to
presentations in which equality reasons are terms of a scoped syntax rather than
lists of elementary steps; the companion article
\cite{ScopedTraceCalculus2026} treats one such presentation.
\end{remark}

\section{A design constraint on any replacement rule set}
\label{sec:design}

Theorem~\ref{thm:rweq-total} is a fact about one rule set.  What must a
\emph{replacement} satisfy to keep the quotient informative?  There is a single
answer, and both the collapse and its sharpness counterpart are instances of it.

Make ``keeping information'' precise.  An \emph{invariant} of a relation \(S\)
on \(X\) is a map \(I:X\to V\) with \(S(x,y)\to I(x)=I(y)\); it is
\emph{nonconstant} when \(I(x)\ne I(y)\) for some \(x,y\).  Invariants of \(S\)
are exactly the functions descending to \(X/S\), so a nonconstant invariant is
precisely a nonconstant function on the quotient.

\begin{theorem}[Totality versus information]
\label{thm:design}
For every setoid \(S\) on \(X\),
\[
 \SetTot(S)
 \quad\Longleftrightarrow\quad
 \text{every invariant of \(S\) is constant.}
\tag{12.1}\label{eq:design}
\]
\end{theorem}

\begin{proof}
Left to right is immediate.  For the converse, apply the hypothesis to the
canonical invariant \(\eta_S:X\to X/S\), which is an invariant by quotient
soundness; constancy gives \(\eta_S(x)=\eta_S(y)\), and quotient exactness
returns \(S(x,y)\).
\end{proof}

Two side conditions deserve explicit mention, since both are easy to drop when
paraphrasing.  First, \eqref{eq:design} needs no inhabitance hypothesis---on the
empty carrier both sides hold vacuously---but the \emph{combination} with
Theorem~\ref{thm:quotient-contractible} does, exactly as
Warning~\ref{warn:empty-carrier} records: for \(X=\varnothing\) every invariant
is vacuously constant while \(X/S\) is empty and not contractible.  Every
statement below is therefore phrased for an inhabited carrier, which the pointed
fibers of interest always are.  Second, the converse direction of
\eqref{eq:design} needs only the single invariant \(\eta_S\), whose codomain
lies in the universe of \(X\); the forward direction holds for invariants valued
in any universe.  So ``every invariant is constant'' may be read unrestrictedly,
the two readings being linked by totality, which mentions no universe at all.

With those recorded, placing Theorem~\ref{thm:design} beside
Theorem~\ref{thm:quotient-contractible} is the point: for an inhabited carrier,
\[
 \isContr(X/S)\iff\SetTot(S)\iff\text{no information survives.}
\tag{12.2}\label{eq:opposition}
\]
The two design goals are therefore not merely difficult to reconcile but
strictly opposed.  A rewrite quotient supports unrestricted based elimination
exactly to the extent that it has stopped distinguishing anything.  This is the
general form of the trade-off that Section~\ref{sec:collapse} exhibited
concretely.

The contrapositive is what one checks a proposal against.

\begin{corollary}[Certificate of nontriviality]
\label{cor:design-certificate}
Exhibiting one invariant that separates two paths proves that the generated
relation is not total, hence that the quotient is not contractible and does not
support unrestricted based elimination.
\end{corollary}

Section~\ref{subsec:sharpness} is exactly such a certificate: trace-length
parity is an invariant of the groupoid fragment separating \(\mathbf r_a\) from
\(\mathbf s_a\), so that fragment is not total.  Presenting the sharpness result
this way shows it is an instance of Theorem~\ref{thm:design} rather than an
unrelated argument.

In the failing direction the mechanism is equally simple, and equally general.

\begin{proposition}[Universal predecessors are fatal]
\label{prop:universal-predecessor}
If some \(e:X\) satisfies \(S(e,x)\) for every \(x\), then \(S\) is total.
Consequently a rule set that keeps any information at a fiber admits no such
\(e\) there.
\end{proposition}

\begin{proof}
Compose the inverse of \(S(e,x)\) with \(S(e,y)\).  The second statement is the
contrapositive, via Corollary~\ref{cor:design-certificate}.
\end{proof}

Lemma~\ref{lem:collapse-rule} says precisely that \(\varepsilon_{a,b}\) is a
universal predecessor, so the collapse is this proposition applied to the trace-free
path.  Nothing about traces, lists, or path algebra is involved; any rule making
one object a one-step predecessor of an entire fiber has the same effect.

We can therefore state the constraint a redesign must meet.

\begin{corollary}[What a redesign must do]
\label{cor:redesign}
Let \(\triangleright'\) be any replacement for the primitive rules and let
\(\approx'\) be its equivalence closure.  If the loop quotient
\(\PathTy_A(a,a)/{\approx'}\) is to be nontrivial, then:
\begin{enumerate}[label=(\roman*),itemsep=1pt,topsep=3pt]
\item \(\approx'\) must admit a nonconstant invariant on \(\PathTy_A(a,a)\);
\item no path may be a \(\triangleright'\)-predecessor of the whole fiber; and
\item consequently \(\PathTy_A(a,a)/{\approx'}\) does \emph{not} support
  unrestricted based elimination, by Theorem~\ref{thm:pathrw-k}.
\end{enumerate}
\end{corollary}

Item (iii) is not a defect of the redesign; it is the price named by
\eqref{eq:opposition}, and it is unavoidable at the level of ordinary setoid
quotients.  An implementation must choose which of the two properties it wants.

Applied to the present rules, the diagnosis is concrete.  The two candidates for
repair are the trace-erasing clause~\eqref{eq:lamcongr-trace} and the
application rule~\eqref{eq:fun-app-beta}.  Giving \(\mathsf{lamCongr}\) a trace,
or restricting~\eqref{eq:fun-app-beta} so that its left-hand side is not
trace-free, removes the universal predecessor;
Proposition~\ref{prop:universal-predecessor} says that is necessary, and
Corollary~\ref{cor:design-certificate} says a proposed replacement is validated
by exhibiting an invariant, exactly as parity validates the groupoid fragment.
We do not carry out the redesign here, but Theorem~\ref{thm:raw-criterion} and
Theorem~\ref{thm:design} between them reduce checking one to two finite tasks:
verify no universal predecessor, and exhibit a separating invariant.

\begin{remark}[Scope]
\label{rem:design-scope}
Theorem~\ref{thm:design} concerns ordinary setoid quotients, as does
Theorem~\ref{thm:quotient-contractible}.  A construction that is not an exact
quotient---a higher quotient, a truncation, or a proof-relevant identity
type---is outside its scope, and is the natural place to look for a target
retaining information while supporting a stronger eliminator.
\end{remark}

\section{Genuine and synthetic circle and torus quotients}
\label{sec:genuine-synthetic}

The repository contains two constructions that earlier names and comments
could invite a reader to conflate:
\begin{enumerate}[label=(\alph*)]
\item the genuine quotient
  \(\PRW_A(a,a)=\PathTy_A(a,a)/\mathsf{RwEq}\); and
\item a separate syntax of winding expressions, quotiented by equality of an
  integer-valued encoding.
\end{enumerate}
They have different carriers, different relations, and, for the present circle
and torus definitions, incompatible contractibility properties.

\subsection{A pointed-subsingleton calculation}

For carriers that are pointed subsingletons there is a second, independent route
to contractibility of the loop quotient, and it is worth recording because it
survives changes to the rule set that Theorem~\ref{thm:rweq-total} does not.

\begin{proposition}[Loop quotient of a pointed subsingleton]
\label{prop:subsingleton-loop-quotient}
If \(\prod_{x:A}(x=a)\), then every raw loop \(p:\PathTy_A(a,a)\) is
\(\mathsf{RwEq}\)-equivalent to \(\mathbf r_a\), and hence
\(\isContr(\PRW_A(a,a))\).
\end{proposition}

\begin{proof}
Every stored source and target is \(a\), and proof irrelevance identifies each
step's equality field with \(\refl_a\).  Induct on the stored list: the empty
list is \(\mathbf r_a\), and otherwise split off the leading singleton, rewrite
it to reflexivity by the transport-reflexivity beta rule, remove the resulting
left unit, and apply the induction hypothesis.  This builds an explicit
composite \(\mathsf{RwEq}\) witness---a rewrite certificate, not an
equality-of-normalizations argument.  Contractibility follows by soundness, or
directly from Theorem~\ref{thm:raw-criterion}, since the first conclusion is
already \(\RwTot(A,a)\).
\end{proof}

Theorem~\ref{thm:rweq-total} subsumes this, dropping the hypothesis entirely.
The proposition is retained because it uses only the transport-reflexivity
erasure and the unit rule: it survives the removal of the function-space rules
contemplated in Section~\ref{sec:design}, whereas the collapse theorem does not.
For pointed subsingletons the loop quotient is thus contractible for two
independent reasons.

\subsection{The current genuine circle and torus}

The artifact currently defines
\[
 \begin{aligned}
  \mathsf{CircleCompPath}&::=\mathsf{base},\\
  \mathsf{Circle}&:=\mathsf{CircleCompPath},\\
  \mathsf{Torus}&:=\mathsf{Circle}\times\mathsf{Circle}.
 \end{aligned}
\tag{13.1}\label{eq:current-circle-torus}
\]
There is no higher-inductive loop constructor in the carrier
\(\mathsf{CircleCompPath}\).  Its separately declared loop-expression syntax
does not alter the equality type of the carrier.

Every element of \(\mathsf{Circle}\) is \(\mathsf{base}\) by ordinary
constructor elimination, and every element of \(\mathsf{Torus}\) is the pair
of base points.  Proposition~\ref{prop:subsingleton-loop-quotient} therefore
gives:

\begin{theorem}[Genuine circle and torus loop quotients]
\label{thm:genuine-circle-torus}
Under the current definitions,
\[
 \begin{aligned}
  \isContr\bigl(
    \PRW_{\mathsf{Circle}}(\mathsf{base},\mathsf{base})
  \bigr),\\
  \isContr\bigl(
    \PRW_{\mathsf{Torus}}(\mathsf{base},\mathsf{base})
  \bigr).
 \end{aligned}
\tag{13.2}\label{eq:genuine-contractible}
\]
Both genuine loop quotients satisfy local quotient-level axiom K and hence the
unrestricted elimination criterion of Theorem~\ref{thm:pathrw-k}.
\end{theorem}

The displayed notation abbreviates the loop fibers
\(\PRW\,\mathsf{Circle}\,\mathsf{base}\,\mathsf{base}\) and similarly for the
torus.  It does not assert that the raw record type is a subsingleton: raw step
lists remain observable before quotienting.

\subsection{The synthetic winding presentations}

Independently, the development defines an indexed syntax of loop expressions
\(e::=\mathsf{loop}\mid\mathsf{refl}\mid\mathsf{symm}(e)\mid
\mathsf{trans}(e,e')\) with the expected integer winding encoding \(w(e)\)
(\(w(\mathsf{loop})=1\), \(w(\mathsf{refl})=0\), \(w\) anti-preserving symmetry
and additive on composition), and forms the setoid quotient
\[
 \mathsf{CircleWind}
 =\mathsf{Expr}/(e\sim e'\Longleftrightarrow w(e)=w(e'))
 \ \Eqv\ \mathbb Z,
\tag{13.3}\label{eq:circle-wind-int}
\]
the equivalence by encode/decode.  The synthetic torus presentation is the
product
\(\mathsf{TorusWind}=\mathsf{CircleWind}\times\mathsf{CircleWind}
\Eqv\mathbb Z\times\mathbb Z\).
Legacy names in the development denote both by \(\pi_1\).  These encode/decode
theorems are valid and useful algebraic presentations; the correction is solely
that they are not the genuine \(\PRW\) loop fibers of the carriers
in~\eqref{eq:current-circle-torus}.

\begin{proposition}[Synthetic noncontractibility]
\label{prop:synthetic-noncontractible}
Both \(\mathsf{CircleWind}\) and \(\mathsf{TorusWind}\) are noncontractible.
\end{proposition}

\begin{proof}
Contractibility is invariant under equivalence, and \(\mathbb Z\) contains
\(0\ne1\) while \(\mathbb Z^2\) contains \((0,0)\ne(1,0)\).
\end{proof}

\subsection{No equivalence bridge}

\begin{theorem}[Circle and torus no-bridge results]
\label{thm:no-bridges}
There is no equivalence between the current genuine circle loop quotient and
\(\mathsf{CircleWind}\), and none between the current genuine torus
loop quotient and \(\mathsf{TorusWind}\):
\[
\begin{aligned}
 \neg\bigl(
   \PRW_{\mathsf{Circle}}(\mathsf{base},\mathsf{base})
   \Eqv\mathsf{CircleWind}
 \bigr),\\
 \neg\bigl(
   \PRW_{\mathsf{Torus}}(\mathsf{base},\mathsf{base})
   \Eqv\mathsf{TorusWind}
 \bigr).
\end{aligned}
\tag{13.4}\label{eq:no-bridges}
\]
\end{theorem}

\begin{proof}
An equivalence transports contractibility.  The source in each line is
contractible by Theorem~\ref{thm:genuine-circle-torus}; the target is
noncontractible by Proposition~\ref{prop:synthetic-noncontractible}.  Either
equivalence yields a contradiction.
\end{proof}

This settles the formerly suggested bridge under the current signatures: it is
not merely unformalized, but impossible.  The exact formal statement excludes
invertible bridges.  It does not claim that no function whatsoever
can be written between the types.

Theorem~\ref{thm:rweq-total} strengthens the conclusion considerably.  The
obstruction is not a defect of the one-constructor carrier: it is present for
every carrier whatsoever.

\begin{theorem}[Universal no-bridge]
\label{thm:universal-no-bridge}
Let \(X\) be any type, \(x:X\), and let \(Y\) be any noncontractible type.  Then
there is no equivalence
\[
 \PRW_X(x,x)\Eqv Y .
\tag{13.5}\label{eq:universal-no-bridge}
\]
In particular \(\PRW_X(x,x)\not\Eqv\mathbb Z\) and
\(\PRW_X(x,x)\not\Eqv\mathbb Z^2\) for every pointed carrier.
\end{theorem}

\begin{proof}
Corollary~\ref{cor:unconditional-repair}(ii) makes the source contractible, and
contractibility transports across an equivalence.
\end{proof}

Theorem~\ref{thm:universal-no-bridge} closes off a natural line of repair.  One
might have hoped that replacing \(\mathsf{CircleCompPath}\) by a genuinely
higher-inductive carrier would make the genuine loop quotient agree with the
synthetic winding presentation.  Under the present rewrite rules that is
impossible for \emph{any} carrier: the genuine loop quotient is contractible
before the carrier is even chosen.  A nontrivial genuine loop quotient therefore
requires changing the rewrite system---subject to the constraint of
Corollary~\ref{cor:redesign}---and not only the point space.

\subsection{Which earlier statements are corrected}
\label{subsec:corrections}

Because this section revises claims that were public in the development, we
state precisely what was wrong and what replaces it.  The corrections concern
naming and interpretation in the artifact and its documentation; no previously
proved theorem is retracted.

\begin{enumerate}[label=(C\arabic*),leftmargin=3.1em]
\item \emph{The \(\pi_1\) aliases.}  Two names in the development, and
  accompanying documentation, presented the winding-expression quotients
  \eqref{eq:circle-wind-int} as fundamental groups of the
  circle and the torus computed by the computational-path machinery.  The
  encode/decode theorems behind~\eqref{eq:circle-wind-int}
  are correct, but they are theorems about a separately introduced syntax of
  loop expressions.  They are not statements about \(\PRW\), and by
  Theorem~\ref{thm:no-bridges} they cannot be transported to it.  The objects
  are now named as synthetic winding presentations, and the phrase ``genuine
  fundamental group'' is reserved for---indeed, replaced by---``genuine
  \(\PRW\) loop quotient''.
\item \emph{The suggested bridge.}  A connection between the genuine and
  synthetic objects was described as unformalized future work.  It is not
  future work: Theorem~\ref{thm:no-bridges} shows no equivalence exists under
  the current definitions, and Theorem~\ref{thm:universal-no-bridge} shows that
  no change of carrier can produce one.
\item \emph{Normalization as a characterization of \(\mathsf{RwEq}\).}  Some
  documentation described the implemented normalization function as
  characterizing rewrite equivalence.  That function returns
  \(\mathsf{stepChain}(\pi_{=}p)\) and therefore depends only on the ambient
  equality; it is a sound invariant, and equality of normal forms is not by
  itself a rewrite certificate.  Derivations in this paper accordingly exhibit
  named primitive rules.  (Theorem~\ref{thm:rweq-total} shows that this
  particular invariant happens to be complete as well, but for a reason that has
  nothing to do with normalization: both sides of the equivalence are
  trivial.)
\item \emph{The status of the local-K hypothesis.}  Earlier drafts of this paper
  left \(\RwTot\) open and stated the \(\PRW\) results conditionally.  The
  conditional statements remain valid, and Section~\ref{sec:collapse} discharges
  the hypothesis.
\end{enumerate}

\begin{table}[t]
\centering
\small
\renewcommand{\arraystretch}{1.18}
\begin{tabular}{>{\raggedright\arraybackslash}p{0.29\textwidth}
                >{\raggedright\arraybackslash}p{0.28\textwidth}
                >{\raggedright\arraybackslash}p{0.31\textwidth}}
\toprule
Object & Defining quotient & Current result\\
\midrule
Genuine circle loop fiber
 & raw loops modulo \(\mathsf{RwEq}\)
 & Contractible; local K holds.\\
Synthetic circle winding quotient
 & expressions modulo equal integer encoding
 & Equivalent to \(\mathbb Z\); noncontractible.\\
Genuine torus loop fiber
 & raw product-carrier loops modulo \(\mathsf{RwEq}\)
 & Contractible; local K holds.\\
Synthetic torus winding quotient
 & product of two synthetic circle quotients
 & Equivalent to \(\mathbb Z^2\); noncontractible.\\
Genuine loop fiber of any pointed carrier
 & raw loops modulo \(\mathsf{RwEq}\)
 & Contractible (Thm.~\ref{thm:rweq-total}); no equivalence to a
   noncontractible type.\\
\bottomrule
\end{tabular}
\caption{The current genuine and synthetic objects have incompatible
contractibility and cannot be equivalent; by the collapse theorem the last row
generalises the first and third to every carrier.}
\label{tab:genuine-synthetic}
\end{table}

\section{Consequences and boundaries}
\label{sec:comparison}

\subsection{Where each stage does its work}

The classification is a chain of equivalences and its value lies in keeping the
stages apart:
\[
 \Jbeta_{R_M}
 \Longleftrightarrow
   \isContr\!\left(\sum_b\sum_hM(b,h)\right)
 \Longleftrightarrow
   \isContr\bigl(M(a,\refl_a)\bigr).
\tag{14.1}\label{eq:new-chain}
\]
The first equivalence is the identity-systems step of
Remark~\ref{rem:standard-stage} and is not specific to metadata; the second is
the coordinate calculation of Theorem~\ref{thm:total-fiber}.  Conflating them
would suggest that every decoration of equality inherits path induction, which is
false.  Quotienting adds a relation-theoretic stage whose second equivalence is
quotient exactness, turning ``collapse enough metadata'' into the exact
requirement ``relate every pair at reflexivity'':
\[
 \Jbeta_{\overline R_S}
 \Longleftrightarrow
 \isContr\bigl(M(a,\refl_a)/S_{a,\refl_a}\bigr)
 \Longleftrightarrow
 \SetTot(S_{a,\refl_a}).
\tag{14.2}\label{eq:repair-chain}
\]
For a generic projection the summary is
\(\Jbeta_R\Longleftrightarrow\isContr(\Ker(q))\); for \(\PRW\) the kernel is the
loop quotient, so the condition becomes local axiom K and then \(\RwTot\), which
Section~\ref{sec:collapse} verifies outright.  Keeping raw data, quotient
classes and ambient equality distinct is what makes that last step informative
rather than circular: what it establishes is the left-hand side, by an argument
entirely about primitive rewrite rules.

\subsection{What the theorems do not say}

\begin{enumerate}[label=(\roman*)]
\item They do not refute ordinary equality elimination: the projection
  \(\proj\) still supports Theorem~\ref{thm:motive-factorization}.
\item They do not say metadata cannot be composed, inverted, normalized, or
  rewritten; those operations are orthogonal to fiber contractibility.
\item Every beta law in the classification theorems is propositional.  The
  eliminator of Theorem~\ref{thm:quotient-path-induction} happens to satisfy a
  judgmental one, for the reason given in Remark~\ref{rem:judgmental-beta}; no
  judgmental classification is claimed.
\item They do not identify the trace-carrying record with an intensional or
  homotopical identity type, model full MLTT or HoTT, or prove anything about
  univalence or higher inductive types.
\item For raw records what is ruled out is one unrestricted based eliminator for
  all motives, not a section for each individual metadata-sensitive motive.
\item Theorems~\ref{thm:universal-repair} and~\ref{thm:design} classify exact
  ordinary setoid quotients; no converse is claimed for higher quotients or for
  quotient-like interfaces without exactness.
\item Theorem~\ref{thm:rweq-total} is a theorem about the primitive rule set
  formalized here, not about computational-path calculi in general
  (Remark~\ref{rem:collapse-scope}).  It is not evidence that \(\PRW\) is a good
  identity-like object---by Corollary~\ref{cor:unconditional-repair}(i) it is a
  reformulation of ambient equality---but Theorem~\ref{thm:design} shows that
  any set-level rewrite quotient supporting unrestricted \(J\) must be equally
  degenerate.
\item The no-bridge theorems concern invertible bridges; they do not claim that
  no function whatsoever can be written between the types.
\end{enumerate}

\subsection{Relation to proof relevance and to computational paths}

Ambient equality here is proof-irrelevant, and the proof-relevance of the record
comes from its \(\U\)-valued metadata field rather than from multiple
inhabitants of an equality proposition: \((L,h)\ne(L',h)\) can hold because
\(L\ne L'\), even though any two witnesses \(h,k:a=b\) are equal.  In
intensional homotopy type theory identity itself may have higher structure and
contractibility is formulated internally \cite{HoTTBook2013}; our results
concern an external refinement of a proof-irrelevant equality and neither
replace the HoTT identity type nor transfer its higher-groupoid interpretation
to a trace record.

Computational-path work treats equality reasons and their rewrites as
first-class syntax
\cite{deQueirozEtAl2016,RamosEtAl2017,RamosEtAl2018,deVerasEtAl2025,
RamosEtAl2026}.  Groupoid laws up to \(\mathsf{RwEq}\) organize that step
calculus; they do not by themselves provide unrestricted dependent elimination
over raw records, and Theorem~\ref{thm:sharpness} shows they are also not what
destroys it.  The reading we recommend is the conjunction of the two halves:
computational paths are a legitimate presentation of equality \emph{reasons},
and their rewrite quotient is a legitimate identity-like object, but under the
present rule set these are different jobs done by different objects, and no
single object does both.

\section{Related work}
\label{sec:related}

Martin-L\"of type theory makes identity elimination and its computation law
foundational \cite{MartinLof1984}.  The groupoid interpretation of
Hofmann and Streicher showed that intensional identity need not collapse to
proof-irrelevant equality \cite{HofmannStreicher1998}.  Homotopical and
weak-factorization semantics further explain identity elimination through
path-object structure \cite{AwodeyWarren2009,GambinoGarner2008}, while weak
higher-groupoid results expose the iterated coherence generated by identity
types \cite{Lumsdaine2010,vanDenBergGarner2011}.

The closest antecedent of our starting point is the theory of \emph{identity
systems}.  A pointed family admits based induction exactly when its total space
is contractible; equivalently, when it is a homotopy-initial algebra for the
associated signature.  This is \cite[Thm.~5.8.2]{HoTTBook2013}, developed from
the initial-algebra semantics of inductive types in
\cite{AwodeyGambinoSojakova2012,AwodeyGambinoSojakova2017}.  We state that
equivalence as Theorem~\ref{thm:j-total} in a proof-irrelevant metatheory with
propositional beta, and we do not claim it as new; a reader who grants it may
start at Section~\ref{sec:metadata}.  What is added here is what happens when the
family is a \emph{decorated} equality.  The total space is then computed
exactly---the endpoint and equality coordinates cancel and precisely one fiber
survives (Theorem~\ref{thm:total-fiber})---the resulting criterion is classified
for all ordinary setoid quotients (Theorem~\ref{thm:universal-repair}), and it is
reformulated so that it can be checked on representatives of a concrete rewriting
system (Theorem~\ref{thm:raw-criterion}) and then decided for that system
(Theorem~\ref{thm:rweq-total}).  The identity-systems theorem says which objects
deserve the name; the results here say which decorations and which quotients
produce such objects, and at what cost.

Two-level type theories carefully separate strict and fibrant equality
\cite{Capriotti2016,AltenkirchCapriottiKraus2016,AnnenkovEtAl2023}.
Our setting is not a two-level type theory: it has ambient equality and a data
record refining it, without a fibrancy judgment.  The methodological
similarity is only the insistence that the available eliminator be attached to
the correct equality notion.

Categorical frameworks for dependent type theory---categories with attributes
and comprehension categories, internal type theory, and their more recent
uniform formulations \cite{Cartmell1986,Dybjer1996,Hofmann1997,Uemura2023}---fix
the structure that an identity type must satisfy in a model.  The present
classification is complementary and deliberately more modest: it takes ambient
equality as given and asks which \emph{refinements} of it retain the based
eliminator.  A refinement failing the criterion is not thereby a defective
implementation; it is simply not an identity type, and it should be exposed
through the projection of Section~\ref{sec:factorization} instead.

The computational-path literature motivates the trace application and the
use of explicit rewrite histories
\cite{deQueirozEtAl2016,RamosEtAl2017,RamosEtAl2018,RamosEtAl2026}.
Earlier work studies a richer label-sensitive calculus, in which equality
reasons are terms of a scoped syntax rather than lists of elementary steps; the
collapse of Section~\ref{sec:collapse} is a statement about the list-based
formalization and does not transfer to such presentations.  Nothing here is a
historical-priority claim about path induction; the results classify a specific
equality-plus-metadata interface and then decide the classification for one
implementation of it.
\section{Formal verification}
\label{sec:verification}

Every mathematical statement above is machine-checked in Lean~4
\cite{deMouraUllrich2021,MetadataJSoftware2026,QuotientPathInductionSoftware2026}.
The development introduces no custom axioms and contains no incomplete proofs.
Lean's axiom report was checked for every displayed theorem: the quotient-based
results depend only on propositional extensionality and quotient soundness, the
projection/kernel criterion and the two purely relational statements of
Section~\ref{sec:design} depend on no axiom at all, and the computed-fiber
results depend only on propositional extensionality.  \emph{Classical.choice} is
not used anywhere, so the classification is constructive relative to those two
principles.

The theorem-by-theorem correspondence with the source modules is deferred to
Appendix~\ref{sec:formalization}, so that no formalization detail interrupts the
mathematics.  The raw de Bruijn calculus developed alongside this work is a
separate, self-contained article \cite{ScopedTraceCalculus2026}.

\section{Conclusion}
\label{sec:conclusion}

For equality decorated by observable metadata,
\[
                 R_M(a,b)=\sum_{h:a=b}M(b,h),
\]
unrestricted based elimination with propositional beta is controlled by one
object: \(M(a,\refl_a)\).  The proof separates a standard universal-property
step from a metadata-specific equivalence:
\[
 \Jbeta_{R_M}
 \Longleftrightarrow
 \isContr\!\left(\sum_bR_M(a,b)\right)
 \Longleftrightarrow
 \isContr\bigl(M(a,\refl_a)\bigr).
\]

For computational traces, empty and singleton reflexive lists are visible and
distinct.  The singleton is already well formed and composable, so the
failure is not caused by malformed chains.  It is caused by noncanonical
reflexive metadata.  Equality-projection motives remain eliminable, giving a
precise and useful boundary rather than a blanket impossibility result.

The diagnosis now has a universal ordinary-quotient repair.  Given setoids
\(S_{b,h}\), the repaired relation supports unrestricted \(J^\beta\) exactly
when \(S_{a,\refl_a}\) is total.  Every successful set-level repair therefore
identifies all reflexivity metadata.  The indiscrete relation is the canonical
instance: maximal as a relation, minimal in retained information, and terminal
for representative-preserving quotient factors.  Trace-length parity shows the
criterion is genuinely about contractibility rather than about the amount
identified: it collapses infinitely many traces and repairs nothing.

The projection/kernel theorem removes the need for a sigma presentation.
For the rewrite quotient's equality projection it yields the exact local-K
criterion rather than a
blanket slogan about rewrite quotienting, and its raw-level form replaces that
criterion by an elementary statement about representatives: every two raw loops
at the base point must admit a rewrite derivation.  An eliminator for the
rewrite quotient is therefore exactly as strong as that connectivity assertion.

That assertion is then decided.  Because the implementation packages pointwise
paths into a path between functions while recording the empty trace, the
primitive application rule makes the trace-free path a one-step predecessor of
every path with the same endpoints.  Hence \(\mathsf{RwEq}\) is total and
\(\PRW_A(a,b)\Eqv(a=b)\): the implemented rewrite quotient is ambient equality,
retains no computational-path information, and admits only constant invariants.
It supports unrestricted based path induction on every carrier, and the
formulation we recommend keeps the reason attached to the fact---\emph{the
implemented rewrite quotient supports path induction because it is equivalent to
ambient equality}.  What has path induction is not a proof-relevant path object
preserving computational reasons; it is a quotient that has stopped
distinguishing them.  The raw record, by contrast, is beyond repair without
quotienting, since its reflexivity fiber is exactly \(\List(A)\) for every
pointed carrier.

Section~\ref{sec:design} makes the trade-off general rather than anecdotal: a
setoid is total exactly when every invariant of it is constant, so on an
inhabited carrier elimination and information are strictly opposed for ordinary
setoid quotients, and a rule set admitting a universal predecessor at a fiber can
retain nothing there.  A parity invariant localises the present failure---the groupoid fragment
does not collapse---so what forces the degeneracy is the function-space rules
together with the trace-erasing packaging convention.  A redesign changing
either one is checked against two finite tasks: exhibit a separating invariant,
and verify that no path is a predecessor of a whole fiber.

Finally, the circle and torus audit separates constructions that legacy names
had obscured.  The genuine loop quotients of the one-constructor circle and its
product torus are contractible; the synthetic winding quotients are equivalent
to \(\mathbb Z\) and \(\mathbb Z^2\) and are noncontractible, so no equivalence
bridges them.  The collapse theorem upgrades this from a fact about two
particular carriers to a fact about all of them: no pointed carrier has a
genuine loop quotient equivalent to a noncontractible type, so no change of
point space can restore the intended reading while the present rules stand.
Histories, costs, certificates, normalization logs, and labels all follow the
same reflexivity-fiber test.  Any stronger claim---a judgmental beta
classification, full MLTT or HoTT, univalence, a new higher-inductive circle, or
a higher-quotient universality theorem---requires additional theory and is
outside the results proved here.

\appendix

\section{The companion Lean artifact}
\label{sec:formalization}

This appendix records the theorem-by-theorem correspondence with the source
modules \cite{MetadataJSoftware2026,QuotientPathInductionSoftware2026}.  The
public source
\path{ComputationalPaths/Path/TypeTheory/MetadataJ.lean} defines
\path{EqualityRefinement}, \path{MetadataTotal},
\path{UnrestrictedBasedEliminator}, and \path{IsContractible}.  Its main
theorems and constructions are:
\begin{itemize}[leftmargin=1.3em,itemsep=1pt,topsep=3pt]
\item \path{unrestricted_based_elimination_iff_contractible}, the equivalence
  \(\Jbeta\Eqv\isContr(\text{total})\);
\item \path{metadataTotalEquiv} and
  \path{contractible_iff_of_equiv}, giving
  \(T_M(a)\Eqv M(a,\refl_a)\) and invariance of contractibility;
\item \path{metadata_fiber_criterion} and
  \path{factorized_motive_eliminator}, formalizing
  Theorems~\ref{thm:metadata-fiber} and~\ref{thm:motive-factorization};
\item \path{pathRefinementEquiv} and
  \path{path_no_unrestricted_based_eliminator}, formalizing
  Theorem~\ref{thm:no-trace-j}; and
\item \path{composable_trace_no_unrestricted_based_eliminator}, formalizing
  Corollary~\ref{cor:composable}.
\end{itemize}
The same module contains a genuine length-two
\(\texttt{Path.trans}\) construction and the nontrivial
\(\texttt{RwEq}\) inverse-cancellation certificate
\texttt{metadataTwoStep\_cancel}.  There are no custom axioms or incomplete
proofs.

The repair module \path{MetadataRepair.lean}, in the same
\path{Path/TypeTheory} directory, contains the new repair and case-study
results:
\begingroup
\footnotesize
\begin{itemize}[leftmargin=1.3em,itemsep=1pt,topsep=3pt]
\item \path{quotient_contractible_iff_nonempty_and_setoidTotal} and
  \path{quotient_contractible_iff_setoidTotal}, using
  \path{Quotient.sound} and \path{Quotient.exact};
\item \path{metadata_quotient_repair_criterion},
  \path{metadata_setoid_total_all_fibers}, and
  \path{metadata_quotient_fibers_contractible};
\item \path{quotient_map_factor_iff_relationIncluded};\\
  \path{factorToIndiscrete_unique}; and\\
  \path{quotientEquivIndiscreteOfTotal};
\item \path{projection_kernel_criterion},
  \path{pathRwQuot_elimination_iff_loop_quotient_contractible}, and
  \path{pathRwQuot_elimination_iff_local_axiomK};
\item the raw-level criterion
  \path{loop_quotient_contractible_iff_rweq_total},
  \path{local_axiomK_iff_rweq_total},\\
  \path{pathRwQuot_elimination_iff_rweq_total},
  \path{elimination_forces_rweq_on_raw_loops}, and the explicit derivation
  \path{doubleSingletonRweqEmpty} with its quotient consequence
  \path{double_singleton_same_pathRwQuot_class};
\item the computed trace fiber \path{stepEquivPoint},
  \path{traceEquivPointList}, \path{path_eq_of_steps_eq},
  \path{pathEquivPointListProd}, \path{loopPathEquivPointList}, and
  \path{raw_loop_fiber_not_contractible};
\item the failing nontrivial repair \path{traceParitySetoid},
  \path{traceParity_identifies_distinct_traces},
  \path{traceParity_not_setoidTotal}, and
  \path{trace_parity_repair_fails};
\item \path{raw_empty_ne_singleton_reflexive_path},
  \path{raw_empty_rweq_singleton_reflexive_path}, and
  \path{raw_empty_singleton_same_pathRwQuot_class}; and
\item \path{genuine_circle_}\allowbreak\path{piOne_contractible} and
  \path{genuine_torus_}\allowbreak\path{piOne_contractible};
  additionally,
  the two synthetic noncontractibility theorems, and
  \path{no_circle_genuine_}\allowbreak\path{synthetic_bridge} and
  \path{no_torus_genuine_}\allowbreak\path{synthetic_bridge}.
\end{itemize}
\endgroup
The pointed-subsingleton proof constructs a multi-step \(\mathsf{RwEq}\)
derivation by induction on the stored trace; the raw singleton theorem uses the
primitive \path{Step.transport_refl_beta} rule.  The axiom footprint is reported
in Section~\ref{sec:verification}.

The results of Section~\ref{sec:collapse} are formalized in a third
module, in the same directory
\cite{QuotientPathInductionSoftware2026}:
\begingroup
\footnotesize
\begin{itemize}[leftmargin=1.3em,itemsep=1pt,topsep=3pt]
\item \path{lamCongr_steps}, recording~\eqref{eq:lamcongr-trace}, and
  \path{stepEmptyTrace}, which is Lemma~\ref{lem:collapse-rule} obtained from
  \path{Step.fun_app_beta} at \path{PUnit};
\item \path{rweqAny}, \path{rweq_total}, \path{rwEqSetoid_total} and
  \path{rwEqTotalOnLoops_always}, formalizing Theorem~\ref{thm:rweq-total};
\item \path{pathRwQuot_subsingleton}, \path{pathRwQuotEquivPLiftEq},
  \path{pathRwQuot_loop_contractible}, \path{pathRwQuot_localAxiomK} and
  \path{pathRwQuot_axiomK}, formalizing
  Corollary~\ref{cor:unconditional-repair};
\item \path{pathRwQuotJ} with \path{pathRwQuotJ_beta}, the unbased form
  \path{pathRwQuotJ'} with \path{pathRwQuotJ'_beta}, the derived
  \path{pathRwQuotTransport}, and the interface version
  \path{pathRwQuotEliminator}, formalizing
  Theorem~\ref{thm:quotient-path-induction};
\item \path{raw_fails_quotient_succeeds}, formalizing
  Theorem~\ref{thm:raw-vs-quotient};
\item \path{no_loop_quotient_equiv_of_not_contractible},
  \path{no_loop_quotient_equiv_int} and
  \path{no_loop_quotient_equiv_int_prod}, formalizing
  Theorem~\ref{thm:universal-no-bridge}; and
\item the sharpness fragment \path{GroupoidStep} with its embedding
  \path{groupoidStepToStep}, its closure \path{GroupoidRwEq}, the invariant
  \path{groupoidStep_traceParity} and \path{groupoidRwEq_traceParity}, and
  \path{groupoid_fragment_not_total}, formalizing
  Proposition~\ref{prop:parity-invariant} and Theorem~\ref{thm:sharpness};
\item the design constraint of Section~\ref{sec:design}:
  \path{SetoidInvariant}, \path{Nonconstant},
  \path{setoidTotal_iff_all_invariants_constant} for
  Theorem~\ref{thm:design}, with
  \path{all_invariants_constant_of_typeU_invariants_constant} extending
  constancy to invariants valued in any universe (Lean cannot quantify over
  universe levels inside a proposition, so the unrestricted reading is routed
  through totality), \path{not_setoidTotal_of_nonconstant_invariant} for
  Corollary~\ref{cor:design-certificate},
  \path{setoidTotal_of_universal_predecessor} with
  \path{no_universal_predecessor_of_nonconstant_invariant} for
  Proposition~\ref{prop:universal-predecessor},
  \path{rwEq_universal_predecessor} and \path{rwEq_invariant_constant} for the
  instance at hand, and
  \path{groupoidSetoid_parity_invariant} with
  \path{groupoidSetoid_not_total_via_design} re-deriving the sharpness result
  from the general theorem; and
\item \path{groupoidSetoid}, \path{groupoidSetoid_not_setoidTotal},
  \path{groupoid_quotient_loop_not_contractible} and
  \path{fragment_fails_full_system_succeeds}, formalizing
  \eqref{eq:fragment-vs-full}.
\end{itemize}
\endgroup
The beta equation \path{pathRwQuotJ_beta} is proved by \path{rfl}, which is the
formal content of Remark~\ref{rem:judgmental-beta}, and the embedding
\path{groupoidStepToStep} is what makes Theorem~\ref{thm:sharpness} a statement
about a genuine sub-system rather than about an unrelated relation.

The release tag \texttt{metadata-quotient-repair-v2} names the stable
theorem-bearing commit cited in~\cite{MetadataJSoftware2026}; the tag is not a
placeholder DOI.  A future Zenodo archive may supplement that Git reference
when a DOI has actually been issued.

The raw scope-indexed de Bruijn development is intentionally not interleaved
with the present theory paper.  It appears as a self-contained companion
article \cite{ScopedTraceCalculus2026}, which develops raw syntax and its
binder algebra, typed substitution, primitive-redex preservation, contextual
and multi-step reduction with a conditional subject-reduction theorem,
syntactic quotients, and source identity programs, in a scope-indexed
de Bruijn presentation \cite{deBruijn1972}.  That article also proves that its
quotient traces are \emph{unlabelled} in a precise sense: evaluation factors
through a label-free syntax that cannot mention a source derivation.

None of those raw-calculus results is needed for
Theorem~\ref{thm:metadata-fiber}; the dependence runs the other way, in that
the companion applies the present fiber criterion to its own trace records.
Keeping the two separate leaves this article a mathematics/type-theory paper
rather than a report on Lean or on a particular de Bruijn implementation.
\begingroup
\small
\bibliographystyle{amsplain}
\bibliography{refs}
\endgroup

\end{document}
